\documentclass[11pt]{article}
\usepackage[T1]{fontenc}
\usepackage[utf8]{inputenc}
\usepackage{amsmath,amsthm,mathtools}
\usepackage{newtxtext,newtxmath}
\usepackage[margin=0.95in,headheight=15pt]{geometry}
\usepackage{microtype}
\usepackage{xcolor,booktabs,array,longtable,enumitem}
\usepackage{ragged2e,tocloft}
\setlength{\cftbeforesecskip}{1pt}
\setlength{\cftbeforetoctitleskip}{0pt}
\setlength{\cftaftertoctitleskip}{10pt}
\usepackage{fancyhdr}
\usepackage{hyperref}
\usepackage[nameinlink,noabbrev]{cleveref}
\definecolor{ink}{HTML}{263D36}
\definecolor{link}{HTML}{385E54}
\hypersetup{colorlinks=true,linkcolor=link,citecolor=link,urlcolor=link,
 pdftitle={Automorphisms of the Surcomplex Numbers},
 pdfauthor={Research article prepared with ChatGPT},
 pdfsubject={Real forms, Hahn symmetries, valuation, topology, and exponential rigidity}}
\pagestyle{fancy}\fancyhf{}
\fancyhead[L]{\small\scshape Automorphisms of the surcomplex numbers}
\fancyhead[R]{\small Real forms and rigidity}
\fancyfoot[C]{\thepage}
\renewcommand{\headrulewidth}{0.3pt}
\setlength{\parskip}{3pt}
\setlength{\parindent}{1.15em}
\setlist{itemsep=2pt,topsep=4pt,leftmargin=1.6em}
\emergencystretch=1.5em
\newtheorem{theorem}{Theorem}[section]
\newtheorem{proposition}[theorem]{Proposition}
\newtheorem{lemma}[theorem]{Lemma}
\newtheorem{corollary}[theorem]{Corollary}
\theoremstyle{definition}
\newtheorem{definition}[theorem]{Definition}
\newtheorem{example}[theorem]{Example}
\theoremstyle{remark}
\newtheorem{remark}[theorem]{Remark}
\newtheorem{question}[theorem]{Question}
\newcommand{\No}{\mathbf{No}}
\newcommand{\On}{\mathbf{On}}
\newcommand{\R}{\mathbb R}
\newcommand{\C}{\mathbb C}
\newcommand{\Q}{\mathbb Q}
\newcommand{\N}{\mathbb N}
\newcommand{\Z}{\mathbb Z}
\newcommand{\K}{\mathbb K}
\newcommand{\OO}{\mathcal O}
\newcommand{\mm}{\mathfrak m}
\newcommand{\U}{\mathcal U}
\newcommand{\id}{\operatorname{id}}
\newcommand{\Aut}{\operatorname{Aut}}
\newcommand{\Hom}{\operatorname{Hom}}
\newcommand{\Fix}{\operatorname{Fix}}
\newcommand{\supp}{\operatorname{supp}}
\newcommand{\lc}{\operatorname{lc}}
\newcommand{\st}{\operatorname{st}}
\newcommand{\re}{\operatorname{Re}}
\newcommand{\im}{\operatorname{Im}}
\newcommand{\rc}{\operatorname{rc}}
\newcommand{\abs}[1]{\lvert#1\rvert}
\newcommand{\angles}[1]{\langle#1\rangle}
\newcommand{\Gh}{\operatorname{Aut}_{\mathrm{ord}}(\Gamma)}
\newcommand{\Gexp}{\operatorname{Aut}_{\exp}(\No)}
\newcommand{\AutL}{\operatorname{Aut}_{L}}
\newcommand{\repoSHA}{\texttt{dcf86662\allowbreak b574988e\allowbreak 75d08d51\allowbreak 5d3a2a5b\allowbreak a7bbc0d7}}
\newcommand{\repo}[1]{\nolinkurl{#1}}% a path in the collection, relative to docs/
\newcommand{\lbl}[1]{\nolinkurl{#1}}% a label of another report, cited by name
\numberwithin{equation}{section}
\begin{document}
\hypersetup{pageanchor=false}
\begin{titlepage}
\centering
\vspace*{0.1cm}
{\large\scshape Single-source research report\par}
\vspace{0.8cm}
{\Huge\bfseries\color{ink} Automorphisms of the\\[0.2em] Surcomplex Numbers\par}
\vspace{0.65cm}
{\Large Real forms, Hahn symmetries,\\ valuation, and rigidity\par}
\vspace{0.7cm}
{\large Research article prepared with ChatGPT\par}
\vspace{0.3cm}
{22 September 2026\par}
\vspace{0.1cm}
{\small Built from one manuscript of batch~20 of the collection,
pinned to repository commit \texttt{dcf8666}\par}
\vspace{0.6cm}
\begin{minipage}{0.93\textwidth}
\begin{abstract}
The phrase ``automorphism of the surcomplex numbers'' has several inequivalent
meanings. For the algebraically closed field $\K=\No(i)$, the automorphisms
fixing $\No$ pointwise are just identity and conjugation. Those preserving
$\No$ setwise form $\Aut(\No)\times C_2$, but this is a proper subgroup of the
unrestricted field automorphisms. We prove this distinction explicitly, using
Hahn-series phase twists that fix every ordinary complex number and every
valuation value while moving the surreal real axis.

We develop the coefficient, character, value-group, and leading-term-kernel
layers of the valued automorphism group; construct automorphisms moving
ordinary real constants by infinitesimals; and separate strong summation from
fine continuity. A further contrast is that nonidentity surcomplex field
automorphisms can be everywhere fine-differentiable with derivative zero.
We examine class-sized homogeneity, fixed fields, involutions, inequivalent
real forms, norm and circle preservation, simplicity, the omega-map, and
exponential structure. The repository's existing exponential faithfulness
argument is incorporated as prior work, not presented as a new discovery.
That repository report classifies the automorphisms preserving the
logarithmic modulus,
$\AutL(\K)\simeq\Gexp\times C_2$; the present report studies plain field
automorphisms, of which those form a proper subgroup of the real-axis
stabilizer.
The final sections identify precise formalization targets and distinguish
proved statements, imported theorems, and unresolved classification problems.
\end{abstract}
\end{minipage}
\vfill
\begin{minipage}{0.93\textwidth}
\small
\textbf{Scope and status.} An AI-assisted research draft, not refereed.
A mathematical research exposition, not a claim of
complete classification or of historical priority. Written proofs are supplied
for the main reductions and constructions. No new Lean verification is
supplied here. Its labels have no dedicated implementation mappings;
existing generic Lean results cover ordered displacement, exponential
rigidity, logarithmic-modulus classification and the valuation kernel
(Section~\ref{saut:sec:position}). The original repository comparison is
pinned to the revision recorded in Section~\ref{saut:sec:repo}; the relations
to the collection as it now stands are in
Section~\ref{saut:sec:position} and Appendix~\ref{saut:app:provenance}.
The word ``automorphism'' is always qualified by the structure it preserves
(Section~\ref{saut:sec:conventions}). Every label carries the prefix
\texttt{saut:}.
\end{minipage}
\end{titlepage}
\hypersetup{pageanchor=true}
\pagenumbering{roman}
{\fontsize{9.5}{11}\selectfont\setlength{\parskip}{0pt}\tableofcontents}
\clearpage
\pagenumbering{arabic}

\section{Which structure is being preserved?}\label{saut:sec:introduction}

Let $\No$ be Conway's real closed ordered field of surreal numbers and put
\[
 \K=\No(i)=\{x+iy:x,y\in\No\},\qquad i^2=-1.
\]
This is the surcomplex field. Its algebraic closedness follows from real
closedness of $\No$. Write
\[
 c(x+iy)=x-iy,
 \qquad N(x+iy)=x^2+y^2,
 \qquad |x+iy|=\sqrt{x^2+y^2}.
\]
The square root is the nonnegative surreal square root. In particular, the
modulus is \emph{surreal-valued}, not an ordinary real-valued norm.
The underlying surreal construction and its normal forms are classical
\cite{Conway,Gonshor}.

The elementary equation $\alpha(i)^2=-1$ forces any field automorphism
$\alpha$ of $\K$ to send $i$ to $i$ or $-i$. It does \emph{not} force
$\alpha(\No)=\No$. This is the first logical distinction on which the subject
turns. Once $\No$ is preserved, order is automatic on that subfield; without
the distinguished subfield, the pure algebraically closed field has no
intrinsic real axis.

Three groups should therefore never be conflated:
\begin{equation}\label{saut:eq:three}
 \Aut(\K/\No)\ \subsetneq\
 G_{\No}:=\{\alpha:\alpha(\No)=\No\}\ \subsetneq\ \Aut(\K).
\end{equation}
The left-hand group fixes $\No$ \emph{pointwise}. The middle group preserves
it \emph{setwise}. The right-hand group has no such condition. We will prove
\[
 \Aut(\K/\No)=\{\id,c\},\qquad
 G_{\No}\simeq\Aut(\No)\times C_2,
\]
and exhibit witnesses to both strict inclusions.

Other reasonable structures produce other answers. One may mark ordinary
$\C$, the natural valuation ring, every individual valuation value, the
Conway monomials, strong Hahn summation, conjugation, modulus, simplicity,
a derivation, or a real exponential. Naming one does not silently name all
the others. For example, preserving every ordinary complex number and every
valuation value still permits a rotation of some real monomials into
imaginary ones.

\subsection{Literature and the repository}

Kuhlmann--Serra developed decompositions of valued Hahn-field automorphisms
\cite{KS}; their work on Hahn groups supplies a second layer of structure
\cite{KSG}. Kaplan--Krapp--Serra apply these ideas to $\No$, including
strong summation and additional exponential structure \cite{KKS}. We use
these as the structural background, while proving the complex-coefficient
reductions and explicit examples below directly.

The supplied repository already contains a manuscript proving faithful
value-group action for surreal exponential automorphisms, with a
logarithmic-modulus application to $\No(i)$ \cite{RepoRigidity}. That argument
is reproduced with attribution in Section~\ref{saut:sec:exponential}. Its result
is important: one must not simply repeat the corresponding question in an
earlier source as unanswered. Section~\ref{saut:sec:audit} separates that
mathematical issue from publication status and priority.

\subsection{Foundations: a collection of class maps}

We work in a class theory such as NBG with global choice. An individual
surreal number is a set, but $\No$, $\K$, and the graphs of their automorphisms
are proper classes. The notation $\Aut(\K)$ denotes a \emph{definable
collection of class maps}; it is not itself a class-object in NBG. A group
isomorphism formula means a unique factorization statement about individual
automorphisms, with the stated composition law. This convention avoids
forming a nonexistent set or class of all class functions.
Likewise, a finite subgroup means a faithful action of an ordinary finite
group $G$, encoded by one class relation in $G\times\K\times\K$ whose
slices are field automorphisms. No set of proper-class graphs is formed;
the fixed field is defined by quantifying over the finite index set $G$.
For a collection of all class automorphisms, a displayed fixed-field
identity is instead read pointwise: it asserts that an element is fixed
by every individual map exactly when it belongs to the named field.
It does not invoke class comprehension with a quantifier over class maps.

Global choice is used for the nonconstructive extension arguments in
Section~\ref{saut:sec:bare}. The explicit Hahn maps need no class-sized
transcendence basis: their values are defined on the set-sized support of
one input at a time. Every infinite sum in this article has a set-sized
index family. No proper-class support is admitted.

For a set-based treatment, most local algebra and all explicit formulas can
instead be read in a Hahn field $\C((t^\Gamma))$ for a set-sized divisible
ordered abelian group $\Gamma$. A self-map requires the chosen parameters to
preserve that workspace. Global homogeneity and the topology of the full
surreal field should not be transferred to an arbitrary fixed workspace.

\subsection{Notation}

Throughout the full-field discussion,
\[
 \Gamma=(\No,+,<),\qquad t^g:=\omega^{-g}.
\]
Thus $t=\omega^{-1}$ is positive infinitesimal. The sign in this convention
is essential. In $t$-notation, supports increase and are well-ordered; in
$\omega$-notation, normal-form exponents decrease.

\begin{center}
\small
\begin{tabular}{@{}>{\RaggedRight\arraybackslash}p{0.22\textwidth}>{\RaggedRight\arraybackslash}p{0.70\textwidth}@{}}
\toprule
Symbol & Meaning\\
\midrule
$\Aut(\K/E)$ & Field automorphisms fixing the indicated subfield $E$ pointwise\\
$G_{\No}$ & Setwise stabilizer of the distinguished surreal real axis\\
$G_v$ & Automorphisms stabilizing the natural valuation ring\\
$v$, $\lc$ & Minimum support exponent and its nonzero coefficient\\
$\U$ & Automorphisms fixing the leading term of every nonzero input\\
$\Gh$ & Ordered \emph{additive-group} automorphisms of $\Gamma$\\
$\st$ & Residue/standard-part map on the finite subring\\
$c$ & The chosen conjugation $x+iy\mapsto x-iy$\\
\bottomrule
\end{tabular}
\end{center}

\subsection{Conventions: one meaning per word}\label{saut:sec:conventions}

This subsection and the next were written when the manuscript joined the
collection (Appendix~\ref{saut:app:provenance}). They fix vocabulary and
state relations. The only mathematics they add is the elementary observation
\eqref{saut:eq:Lchain} below, which is marked as an observation of the edit.

\paragraph{Kinds of automorphism.} An unqualified ``automorphism'' of $\K$
or of $\No$, and the notation $\Aut(\K)$, $\Aut(\No)$, always mean a
\emph{plain field automorphism}: a bijective ring map, with no further
condition. Every other kind is named by the structure it preserves.
\begin{itemize}
\item A \emph{valued automorphism} is a field automorphism $\alpha$ with
$\alpha(\OO)=\OO$, that is, an element of the group $G_v$ of
Section~\ref{saut:sec:decomposition}. Equivalently
$v(\alpha z)=\tau_\alpha(v z)$ for an ordered automorphism $\tau_\alpha$ of
the value group. A valued automorphism need not fix any value. Every field
automorphism of $\No$ is valued (Lemma~\ref{saut:lem:order} and
Section~\ref{saut:sec:decomposition}); a field automorphism of $\K$ need not
be (Proposition~\ref{saut:prop:nondefinability}).
\item A valued automorphism is \emph{value-fixing} when $\tau_\alpha=\id$,
that is, $v(\alpha z)=v(z)$ for every $z\ne0$. This is the meaning of
``exactly value-preserving'' and of ``values fixed'' in
Theorem~\ref{saut:thm:phase} and Appendix~\ref{saut:sec:examplesummary}.
\item A \emph{$1$-automorphism} is an element of $\U$ in \eqref{saut:eq:U}
(or of $\U_{\R}$ for $\No$): $v(\alpha z-z)>v(z)$ for every $z\ne0$, which
is the same as fixing every leading term. It is value-fixing, and it fixes
$i$. The phase twists $P_\theta$, $\theta\ne0$, are value-fixing but are not
$1$-automorphisms. This leading-term notion is the one of
Kaplan--Krapp--Serra \cite{KKS} that the collection's rigidity report
checked against the pinned version (its Corollary~5.2 uses only that a
$1$-automorphism fixes every value).
\item An \emph{exponential automorphism} is a field automorphism $\sigma$ of
$\No$ with $\sigma(\exp x)=\exp(\sigma x)$ for all $x\in\No$; they form
$\Gexp$. The notion is defined for $\No$ only: no surcomplex exponential is
used in this report. By \eqref{saut:eq:expfaithful}, the identity is the
only exponential automorphism that is value-fixing, and a fortiori the only
exponential $1$-automorphism.
\item An \emph{$L$-automorphism} is a field automorphism $\alpha$ of $\K$
with $\alpha(L(z))=L(\alpha z)$ for every $z\ne0$, where $L(z)=\log|z|$
\eqref{saut:eq:L}. They form $\AutL(\K)$, the notation of the rigidity report;
the source wrote $\Aut(\K,L)$.
\item \emph{Strong} and \emph{strongly $\C$-linear} are the support
conditions of Section~\ref{saut:sec:hahn}. They are not continuity
conditions (Section~\ref{saut:sec:topology}).
\end{itemize}
The inclusions that hold are
\[
 \U\subseteq\{\text{value-fixing}\}\subseteq G_v\subseteq\Aut(\K),
 \qquad
 \{\id,c\}=\Aut(\K/\No)\subseteq\AutL(\K)\subseteq G_{\No}\subseteq G_v .
\]
The first chain is immediate from the definitions. In the second, the
middle inclusion is the range argument of Theorem~\ref{saut:thm:L} and the
last is Proposition~\ref{saut:prop:continuous}.
\emph{Observation of the edit.} Combining Theorem~\ref{saut:thm:axis} with
Theorem~\ref{saut:thm:L},
\begin{equation}\label{saut:eq:Lchain}
 \AutL(\K)=\{\alpha\in G_{\No}:\ \alpha|_{\No}\in\Gexp\}
 \ \subsetneq\ G_{\No}.
\end{equation}
The equality is Theorem~\ref{saut:thm:L} read through the factorization
\eqref{saut:eq:lift}. Strictness is witnessed by $S_q$ for any rational
$q>0$, $q\ne1$: it lies in $G_{\No}$ (Section~\ref{saut:sec:explicit}), and
its restriction to $\No$ induces the value-group dilation $g\mapsto qg$,
which no exponential automorphism induces
(Proposition~\ref{saut:prop:rationalnonlift}). Whether the first inclusion
$\{\id,c\}\subseteq\AutL(\K)$ is strict is the question whether $\Gexp$ is
nontrivial; neither this report nor the rigidity report decides it.

\paragraph{Rigidity.} In this report the word labels three specific
statements: exact norm rigidity (Theorem~\ref{saut:thm:norm}), the
simplicity statement of Section~\ref{saut:sec:exponential}, and the
faithful value-group action of exponential automorphisms
(Theorem~\ref{saut:thm:expfaithful}); the last is faithfulness, not
triviality, exactly as in the rigidity report. None of them is the
all-scale polynomial rigidity \lbl{f:thm-rigidity} of
\repo{surcomplex/analysis} or the scalar rigidity of
\repo{surreal/gamma-functions}.

\paragraph{Phase.} The phases of the twists $P_\theta$ are the coefficient
multipliers $e^{i\theta\ell(g)}$ of \eqref{saut:eq:phase}. They are not an
argument of $z$, not the finite-angle $\operatorname{cis}$ or the global
phase of \repo{surcomplex/trigonometry}, and not the phase of the
differential-equations report; \repo{NOTATION.md} lists those meanings.

\paragraph{Real form.} A real form is a real closed subfield $F\subset\K$
with $\K=F(i)$ (Section~\ref{saut:sec:realforms}). In
\lbl{ent:prop:real-form} of \repo{surcomplex/entire-functions-at-arbitrary-rank}
and in \lbl{analytic:subsec:fixed-real} of \repo{surcomplex/analytic-geometry}
the phrase names the real-coefficient version of a theorem, an unrelated use.

\paragraph{Symbols.} Here $\Gamma=(\No,+,<)$ is a proper class; this
departs from the recommendation of \repo{NOTATION.md} that $\Gamma$ be a
set-sized group, and the set-sized reading is the one described in the
previous subsection. The conventions $t^g=\omega^{-g}$, $v=\min\supp$ and
$v(0)=\infty$ are the collection's. $\K=\No(i)$ is the field that
\repo{NOTATION.md} writes $\mathrm{SC}=\No[i]$ and the analysis report
writes $\K$; $K_1$, $K_2$ and $F(i)$ in Sections~\ref{saut:sec:axis}
and~\ref{saut:sec:bare} are generic. Six symbols of the source were
renamed so that each has one meaning.
\begin{itemize}
\item The source wrote $E_{\rho,\tau}$ for the coefficient-and-exponent lift
and $E_d$ for the Taylor motion of Theorem~\ref{saut:thm:coeffflow}. Here
they are $M_{\rho,\tau}=M_{\rho,\tau,1}$ and $\Phi_d$, so that $E$ denotes
only an ordered exponential (Theorem~\ref{saut:thm:expfaithful}), as in the
rigidity report.
\item The derivation of Theorem~\ref{saut:thm:shiftflow} is $\mathcal D$
(source: $D$). A $D$ with a subscript is always a diagonal automorphism
$D_\lambda$ or $D_\chi$; the unsubscripted $D(x)=\sigma(x)-x$ in
Lemma~\ref{saut:lem:displacement} is the displacement, as in the rigidity
report.
\item The real Puiseux field of Section~\ref{saut:sec:topology} is
$\mathcal P$ (source: $P$), leaving $P_\theta$ for the phase twists.
\item The automorphism in the non-density argument at the end of
Section~\ref{saut:sec:topology} is $\psi$ (source: $\tau$), since $\tau$
always denotes an ordered automorphism of the value group.
\item The sides of a cut are $(\mathcal L,\mathcal R)$ (source: $(L,R)$),
since $L$ is the logarithmic modulus.
\end{itemize}
$T_s$ denotes the flows of Theorem~\ref{saut:thm:shiftflow}; the maps $T_r$
and $T_{s_0,c}$ of the rigidity report are value-group maps, and for
rational $q>0$ the dilation $S_q$ of \eqref{saut:eq:dilation}, restricted to
$\No$, is that report's canonical field lift $\widehat{T_q}$
(\lbl{prop:hahn-lift}). The coefficient $\ell(g)=[g]_0$ of
\eqref{saut:eq:ell} is that report's layer coefficient $[\gamma]_{s_0}$ at
$s_0=0$. The letter $H$ is local to three places: the kernel subgroup in
the proof of Theorem~\ref{saut:thm:expfaithful}, the real closure in
Theorem~\ref{saut:thm:inequivalentrealform}, and the half-plane $H_F$ of
Section~\ref{saut:sec:axis}.

\subsection{Relation to other reports of the collection}\label{saut:sec:position}

The source compared itself with the collection at its pin; that comparison
is kept in Section~\ref{saut:sec:repo}. The relations below were checked
against the collection as it stands after this report was placed.

\paragraph{Exponential automorphism rigidity.}
\repo{surreal/exponential-automorphism-rigidity} studies exponential
automorphisms of an ordered field with a total exponential and, in its
Section~9, $L$-automorphisms of $F(i)$ for real closed $F$; this report
studies plain field automorphisms of $\K$. The two meet in
Section~\ref{saut:sec:exponential}, whose results are that report's and
are printed here once, with its proofs.
\begin{itemize}
\item Lemma~\ref{saut:lem:displacement} is its \lbl{lem:amplify}
(Lemma~3.1), with the same probe.
\item Theorem~\ref{saut:thm:expfaithful} is the faithfulness of its
\lbl{thm:cofinal} and \lbl{cor:faithful} (Theorem~4.1, Corollary~4.2), by
the route of \lbl{lem:bridge} and \lbl{lem:bounded}; the identification of
the kernel with the finite surreals after it is \lbl{prop:finite-log}.
That report proves more: cofinal and coinitial
displacement, coarsenings, a version needing one visible exponential scale,
and an order-free detector.
\item The display \eqref{saut:eq:expfaithful} and the sentence after it are
its \lbl{eq:main-no}, \lbl{thm:no} and \lbl{cor:question54}
(Theorem~5.1, Corollary~5.2): every
exponential $1$-automorphism of $\No$ is the identity, the negative answer
to Question~5.4 of \cite{KKS}, arXiv version~3. The answer is that report's,
not this one's, and the reading of the question rests on its check of the
pinned PDF.
\item Proposition~\ref{saut:prop:rationalnonlift} is its \lbl{thm:dilation}
(Theorem~8.4), by the second, multiplication-based of its two proofs, with
an arbitrary positive infinite $X$ in place of $\omega$.
\item Theorem~\ref{saut:thm:L} is its \lbl{thm:L-classification} with
\lbl{lem:range} and the kernel clause of \lbl{thm:complex-kernel} and
\lbl{cor:surcomplex} (Lemma~9.1, Theorems~9.2 and~9.4, Corollary~9.5),
specialized to $F=\No$. That report proves it for every real closed $F$
with an ordered exponential, and on every invariant nontrivial convex value
quotient.
\end{itemize}
What this report adds on the surcomplex side is the plain-field picture
around that classification. Theorem~\ref{saut:thm:axis} is the field-level
counterpart: $G_{\No}\simeq\Aut(\No)\times C_2$, with $\Aut(\No)$ in place of
$\Gexp$. That report remarks that a shorter route to its classification,
assuming commutation with conjugation and preservation of the exponential on
the real axis from the start, proves strictly less;
Theorem~\ref{saut:thm:axis} is the exponential-free first step of that
shorter route. By \eqref{saut:eq:Lchain}, $\AutL(\K)$ is a proper subgroup of
$G_{\No}$, which in turn is a proper subgroup of $G_v$ and of $\Aut(\K)$
\eqref{saut:eq:strictchain}. By the kernel clause of
Theorem~\ref{saut:thm:L}, an $L$-automorphism that fixes every value is
$\id$ or $c$. The explicit maps $D_\lambda$, $P_\theta$, $\Phi_d$ and $T_s$
are value-fixing and fix $i$, so a nonidentity one is never an
$L$-automorphism; nor is $S_q$ for rational $q\ne1$, by
\eqref{saut:eq:Lchain}. For other $a$ this report does not decide whether
$S_a|_{\No}$ is exponential. The open
question on the image of exponential automorphisms in the value group
(Section~\ref{saut:sec:questions}) is that report's
\lbl{q:image} (Question~13.1).

\paragraph{Foundations and analysis.} The pair construction of $\K$ and its
algebraic closedness are \lbl{found:prop:complex} of
\repo{foundations-and-computation/foundations}.
Proposition~\ref{saut:prop:discrete} is \lbl{found:thm:discrete} there,
items~(1)--(2), with the absence of isolated points; \repo{FORMALIZATION.md}
records Lean proofs of those clauses for the actual surcomplex carrier, with
lower-universe smallness. The fine derivative of
Section~\ref{saut:sec:derivative} is \lbl{found:eq:derivative}.
The collection's earlier zero-derivative witnesses are locally constant:
the indicator of the finite ring in \lbl{found:rem:fourtypes} and in
\lbl{b:L} of \repo{surcomplex/analysis}, with a Lean proof recorded for the
former. The dilation $S_2$ of Theorem~\ref{saut:thm:zeroderiv} is a
different kind of witness: it is injective, so it is constant on no open
set with more than one point, and the full field has no isolated points.

\paragraph{Diagonal character twists.} The automorphisms
$D_\chi=M_{\id,\id,\chi}$ are the character twists that the collection
already uses on set-sized Hahn fields as tools:
\lbl{spec:lem:finiteindex} of \repo{surcomplex/spectral-theory}, whose
Galois group $\Hom(\Lambda/\Gamma,\mu_q)$ acts by them;
\lbl{hnd:lem:supportdegree} of
\repo{surcomplex/hidden-negative-hermitian-directions}; and the twists
built from \lbl{bst:lem:characters} in
\repo{surreal/transcendence-over-bounded-support}. There the characters
take roots of unity or other values on a set-sized group; here $\chi$ is
an arbitrary homomorphism $\Gamma=\No\to\C^\times$, and the twists are
placed in the decomposition of Theorem~\ref{saut:thm:decomp}.

\paragraph{Unit circle and rotations.} The Cayley map in the proof of
Theorem~\ref{saut:thm:circle}, $u=(x-i)/(x+i)$, is a reparametrization of
\lbl{trigonometry:thm:cayley} of \repo{surcomplex/trigonometry}
($u=C(-1/x)$ for $x\ne0$, while $x=0$ gives $u=-1$);
that report identifies the unit circle with
$\operatorname{SO}(2,\No)$ in Section~\lbl{trigonometry:sec:circle}.
Rotations there, and in \repo{surreal/euclidean-three-space}, placed in the
same batch, are $\No$-linear isometries, not field automorphisms. The
contrast is exact: multiplication by a unit-circle element preserves every
modulus, whereas by Theorem~\ref{saut:thm:norm} the only field
automorphisms preserving every modulus are $\id$ and $c$.

\paragraph{Lean interfaces.} Besides the two modules the source inspected
(Section~\ref{saut:sec:repo}), the collection has
\nolinkurl{Surreal/Surcomplex/Basic.lean}, whose \texttt{conj} is conjugation
as a ring automorphism of the concrete surcomplex field, and
\nolinkurl{Surreal/HahnSeries/ComplexNumbers.lean}, whose
\texttt{coefficientEquiv} lifts a ring equivalence of coefficients to Hahn
series over an ordered exponent monoid, keeping supports: the case
$M_{\rho,\id,1}$ of \eqref{saut:eq:monomial} in a workspace. \nolinkurl{Surreal/Algebra/SigmaDerivation.lean}
proves the valued-field forms of the rigidity report's displacement results.
The later \nolinkurl{Surreal/Algebra/ExponentialProfile.lean} proves its
ordered displacement lemma (\texttt{displacement\_cofinal}) and the
value-fixing exponential rigidity implication (\texttt{eq\_id\_of\_commute}).
The modules \nolinkurl{Surreal/Algebra/LogModulusClassification.lean} and
\nolinkurl{Surreal/Algebra/ComplexValuationKernel.lean} now also prove the
generic logarithmic-modulus classification and valuation-kernel theorem.
They work over an ordered field with nonnegative square roots and an
injective \texttt{OrderedExp}; injectivity is an explicit hypothesis,
not a field of that structure. The kernel theorem additionally uses a
nontrivial convex valuation. The classification includes the direct-product
isomorphism (\texttt{autLEquiv}); the kernel module uses separate definitions
of \texttt{log}, $L$ and the $L$-automorphism predicate, without a formal
bridge to the classification module. Its coarsened action is expressed
relationally, not through a constructed ordered quotient.
These generic results cover the corresponding arguments here;
their mappings use the rigidity report's labels. No dedicated
\texttt{saut:} implementation mapping is added. Instantiation at the actual
surreal exponential, rational non-lifting and the classification for the
actual surcomplex logarithmic modulus remain pending; no class automorphism-group object follows
merely from a theorem about individual type-based maps.

\section{Conjugation and the real-axis stabilizer}\label{saut:sec:axis}

The results of this section apply to any real closed ordered field $F$,
with $K=F(i)$. No infinite summation or specifically surreal theorem is needed.

\begin{lemma}\label{saut:lem:order}
Every field automorphism of a real closed field is order preserving. It fixes
its real algebraic numbers pointwise.
\end{lemma}
\begin{proof}
In a real closed field, $x>0$ is equivalent to $x=y^2$ for some $y\ne0$.
A field automorphism preserves this condition and hence preserves order.
It fixes $\Q$. Each real algebraic number is isolated from the other real
roots of its minimal polynomial by a rational interval, so it is fixed too.
\end{proof}

\begin{theorem}[Centralizer and stabilizer]\label{saut:thm:axis}
For $K=F(i)$ with $F$ real closed,
\begin{align}
 \Aut(K/F)&=\{\id,c\},\label{saut:eq:overF}\\
 \{\alpha\in\Aut(K):\alpha(F)=F\}
 &=C_{\Aut(K)}(c)=N_{\Aut(K)}(\angles{c})\nonumber\\
 &\simeq\Aut(F)\times C_2.\label{saut:eq:axis}
\end{align}
The pair $(\sigma,e)$, with $e\in\{0,1\}$, acts by
\begin{equation}\label{saut:eq:lift}
 (x+iy)\longmapsto\sigma(x)+(-1)^e i\sigma(y).
\end{equation}
\end{theorem}
\begin{proof}
The roots of $X^2+1$ in $K$ are exactly $i,-i$. Thus an automorphism fixing
$F$ has precisely the two possibilities in \eqref{saut:eq:overF}. If $\alpha$
preserves $F$, its restriction is a field automorphism $\sigma$ and its
value on $i$ determines \eqref{saut:eq:lift}. Conversely, this formula respects
addition and the multiplication rule
\[
 (x+iy)(u+iv)=(xu-yv)+i(xv+yu),
\]
and has an inverse of the same form. The lift $\widehat\sigma$ fixing $i$
commutes with $c$, so the product is direct.

If $\alpha$ commutes with $c$, it carries the fixed field $\Fix(c)=F$ onto
itself. The converse follows from \eqref{saut:eq:lift}. Finally, normalizing the
order-two subgroup $\{\id,c\}$ means fixing its unique nonidentity member
under conjugation, which is the centralizer condition.
\end{proof}

For surcomplex numbers, $\Aut(\No)$ is very far from trivial. In
Section~\ref{saut:sec:explicit}, exponent dilation will give a real-axis-preserving
automorphism sending $\omega$ to $\omega^2$. Thus even retaining conjugation
leaves much more than the two classical-looking maps.

\subsection{Equivariance is not exact preservation}

For $\alpha=\widehat\sigma c^e$ one has
\begin{equation}\label{saut:eq:equivariantnorm}
 N(\alpha z)=\sigma(Nz),\qquad |\alpha z|=\sigma(|z|).
\end{equation}
The second identity follows because $\sigma$ preserves nonnegativity and
square roots. In particular, balls are taken to balls, but their radii may
change. Calling every such map an isometry would be misleading: it is
\emph{semilinearly} compatible with the surreal-valued modulus.

\begin{theorem}[Exact norm rigidity]\label{saut:thm:norm}
A field automorphism $\alpha$ of $F(i)$ satisfying
$N(\alpha z)=N(z)$ for every $z$ is either $\id$ or $c$.
The same conclusion holds under exact modulus preservation.
No real-axis-preservation hypothesis is needed.
\end{theorem}
\begin{proof}
The identity
\[
 N(z+1)-N(z)-1=2\re z
\]
and $\alpha(z+1)=\alpha z+1$ imply $\re(\alpha z)=\re z$.
For $x\in F$, write $\alpha x=x+iy$. Exact norm preservation gives
$x^2+y^2=x^2$, hence $y=0$. Thus $\alpha$ fixes $F$ pointwise and
\cref{saut:thm:axis} applies. Exact modulus preservation implies exact norm
preservation by squaring.
\end{proof}

\subsection{The unit circle remembers the real axis}

Let $U_F=\{z\in F(i):N(z)=1\}$. Merely preserving this set, without preserving
individual norms, already recovers the real-axis stabilizer.

\begin{theorem}\label{saut:thm:circle}
For a field automorphism $\alpha$ of $F(i)$,
\[
 \alpha(U_F)=U_F\quad\Longleftrightarrow\quad\alpha(F)=F.
\]
\end{theorem}
\begin{proof}
The reverse implication follows from \eqref{saut:eq:equivariantnorm}. For the
forward implication use the Cayley parametrization
\[
 u=\frac{x-i}{x+i},\qquad
 x=i\frac{1+u}{1-u}.
\]
For every $x\in F$, the first expression lies in $U_F\setminus\{1\}$.
For every $u\in U_F\setminus\{1\}$, the second expression is real: conjugating
it and using $\bar u=u^{-1}$ leaves it unchanged. Applying $\alpha$ to the
inverse expression changes $i$ at most to $-i$, so $\alpha x\in F$.
Apply the same reasoning to $\alpha^{-1}$ to obtain equality of the axes.
\end{proof}

Let $H_F=\{z:\im z>0\}$ be the chosen upper half-plane. If a field
automorphism preserves $H_F$ setwise, it also preserves $-H_F$, so it
preserves their complement $F$. Since $i\in H_F$, it must fix $i$.
Conversely, a lift $\widehat\sigma$ preserves $H_F$ by order preservation.
Thus its stabilizer consists exactly of the lifts $\widehat\sigma$. There is no compatible ordered-field structure on $F(i)$
itself, since $i^2=-1$.

\section{Normal forms, valuation, and summability}\label{saut:sec:hahn}

Conway normal form identifies the ordered field $\No$ with real Hahn series
having surreal exponents and set-sized well-ordered support. Taking real and
imaginary parts coefficientwise gives
\begin{equation}\label{saut:eq:hahn}
 \No=\R((t^\Gamma)),\qquad
 \K=\C((t^\Gamma)),\qquad \Gamma=(\No,+,<).
\end{equation}
The displayed equalities mean canonical identifications. A surcomplex normal
form is
\[
 z=\sum_{g\in S}z_g t^g,
 \qquad z_g\in\C,\quad S\subseteq\Gamma\text{ a well-ordered set}.
\]
For $z\ne0$, define
\[
 v(z)=\min\supp(z),\qquad \lc(z)=z_{v(z)},\qquad
 \operatorname{lt}(z)=\lc(z)t^{v(z)},
\]
and set $v(0)=\infty$. Multiplication is convolution, with only finitely many
contributions to any fixed exponent. Its basic laws are
\[
 v(zw)=v(z)+v(w),\qquad
 v(z+w)\ge\min\{v(z),v(w)\}.
\]
Conjugation acts coefficientwise. Since
$\lc(z\bar z)=|\lc(z)|_{\C}^{\,2}>0$,
\begin{equation}\label{saut:eq:normvalue}
 v(\bar z)=v(z),\qquad v(Nz)=2v(z),\qquad v(|z|)=v(z).
\end{equation}

The valuation ring and its maximal ideal are
\[
 \OO=\{z:v(z)\ge0\},\qquad \mm=\{z:v(z)>0\}.
\]
They are exactly the finite surcomplex numbers and the infinitesimals,
respectively. Every $z\in\OO$ has a unique decomposition $z=a+\varepsilon$
with $a\in\C$, $\varepsilon\in\mm$, and $\st(z)=a=z_0$. Thus
$\OO/\mm\simeq\C$. The analogous real residue field is $\R$.

\begin{definition}[Strong summation]
A set-indexed family $(z_j)_{j\in J}$ is Hahn-summable if the union of its
supports is well-ordered and every exponent occurs in only finitely many
of the supports. Its sum is then defined coefficientwise by finite sums.
An additive map is \emph{strongly additive} when it carries summable
families to summable families and commutes with their sums. An automorphism
is \emph{strong} here when it and its inverse are strongly additive.
A strongly $\C$-linear map is a strongly additive $\C$-linear map;
for an automorphism we also require its inverse to be strongly additive.
For field automorphisms, $\C$-linearity is equivalent to fixing $\C$ pointwise.
\end{definition}

This definition is a support condition, not a limit of finite partial sums.
For example, $\sum_{n\ge0}t^n$ is a legitimate Hahn sum and equals
$(1-t)^{-1}$ even in a setting where its partial sums do not converge in the
full fine topology. Section~\ref{saut:sec:topology} explains that distinction.

\begin{lemma}[Support control used below]\label{saut:lem:support}
If $S$ is a well-ordered subset of an ordered abelian group and $\delta>0$,
then $S+\N\delta$ is well-ordered, and each exponent has only finitely many
representations $s+n\delta$ with $s\in S$ and $n\in\N$.
More generally, the sum of two well-ordered subsets is well-ordered and has
finitely many decompositions of any fixed element.
\end{lemma}
\begin{proof}
Every infinite sequence in a well-order has an infinite nondecreasing
subsequence. A strictly decreasing sequence of sums, by successively taking
such subsequences for its two coordinates, would have a nondecreasing
subsequence, a contradiction. Thus the sum is well-ordered. Infinitely many
distinct representations of one value would give infinitely many distinct
first coordinates. Choosing a strictly increasing subsequence of these
forces a strictly decreasing sequence of second coordinates, again
impossible. Apply this to $S$ and $\N\delta$.
\end{proof}

These finite-decomposition arguments justify the rearrangements in the
constructions below. More general support-controlled formal exponentials
and derivations are developed in \cite{BKKS}; we will need only explicit
fixed-shift cases, which can be checked directly.

\section{Explicit monomial automorphisms and phase twists}\label{saut:sec:explicit}

Let $\rho\in\Aut(\C)$, let $\tau\in\Gh$, and let
$\chi:(\Gamma,+)\to(\C^\times,\cdot)$ be a group homomorphism. Define
\begin{equation}\label{saut:eq:monomial}
 M_{\rho,\tau,\chi}\left(\sum_g z_g t^g\right)
   =\sum_g\rho(z_g)\chi(g)t^{\tau(g)}.
\end{equation}
Here $\chi$ need not arise from an exponential of an additive complex-valued
map. Arbitrary multiplicative characters are allowed.

\begin{theorem}[Monomial construction]\label{saut:thm:monomial}
Formula~\eqref{saut:eq:monomial} defines a strong field automorphism of $\K$, with
\[
 v(Mz)=\tau(vz),\qquad
 \lc(Mz)=\rho(\lc z)\chi(vz).
\]
For composition in the order $M_1M_2$, the parameters are
\begin{equation}\label{saut:eq:monomiallaw}
 \rho=\rho_1\rho_2,\quad \tau=\tau_1\tau_2,\quad
 \chi(g)=\rho_1(\chi_2(g))\chi_1(\tau_2(g)).
\end{equation}
The map preserves the real axis exactly when
\begin{equation}\label{saut:eq:monomialreal}
 \rho\in\{\id_{\C},c|_{\C}\},\qquad
 \chi(g)\in\R_{>0}\quad\text{for every }g.
\end{equation}
\end{theorem}
\begin{proof}
The order bijection $\tau$ takes a well-ordered support to a well-ordered
support. All transformed coefficients on the support are nonzero.
Additivity is coefficientwise. For multiplication, each coefficient is a
finite sum, $\tau(g+h)=\tau(g)+\tau(h)$, and
$\chi(g+h)=\chi(g)\chi(h)$, so convolution is preserved. The same formula
with inverse parameters gives a two-sided inverse; explicitly one can use
$\rho^{-1},\tau^{-1}$ and
$g\mapsto\rho^{-1}(\chi(\tau^{-1}g)^{-1})$. Support unions and their finite
incidence multiplicities are preserved, proving strongness in both
directions. Substitution proves \eqref{saut:eq:monomiallaw}.

If the real axis is preserved, the image of each real constant is a real
constant, so $\rho(\R)=\R$. A field automorphism of $\R$ is the identity,
because order and rational cuts determine each real number. Consequently
$\rho$ is identity or conjugation on $\C$. Also $M(t^g)$ is real, so
$\chi(g)\in\R^\times$. Divisibility of $\Gamma$ gives
$\chi(g)=\chi(g/2)^2>0$. Conversely, \eqref{saut:eq:monomialreal} makes the image
of every real-coefficient series real, and the same holds for the inverse.
\end{proof}

\subsection{Rescaling all valuation exponents}

For every positive surreal $a$ let
\begin{equation}\label{saut:eq:dilation}
 S_a\left(\sum_g z_g t^g\right)=\sum_g z_g t^{ag}.
\end{equation}
Multiplication by $a$ is an ordered additive automorphism of $\Gamma$, so
\cref{saut:thm:monomial} applies. We obtain
\[
 S_aS_b=S_{ab},\qquad S_a^{-1}=S_{1/a},\qquad
 S_a(t)=t^a,\qquad S_a(\omega)=\omega^a.
\]
Every $S_a$ fixes ordinary $\C$ pointwise, commutes with $c$, and preserves
order on $\No$. It fixes individual valuation values only for $a=1$.
For instance,
\[
 S_2\left(3+2t^{-1}+(1+i)t^{1/2}\right)
 =3+2t^{-2}+(1+i)t.
\]
This proves the first strict inclusion in \eqref{saut:eq:three}.

Notice that $S_a$ generally does \emph{not} commute with the omega-map.
Indeed, $S_a(1)=1$ whereas
$S_a(\omega^1)=\omega^a\ne\omega^{S_a(1)}$ when $a\ne1$.
The omega-map must not be mistaken for an exponent reindexing convention.

\subsection{Positive real characters}

For an additive map $\lambda:\Gamma\to\R$, define
\[
 D_\lambda\left(\sum_g z_g t^g\right)
 =\sum_g z_g e^{\lambda(g)}t^g.
\]
These are strongly $\C$-linear real-axis-preserving automorphisms fixing
all valuation values. Their group law is
$D_\lambda D_\mu=D_{\lambda+\mu}$. They generally change leading
coefficients, so they are not leading-term-fixing automorphisms.

There are explicit nonzero $\lambda$. Regard an exponent $g$ itself as a
surreal normal form and set
\begin{equation}\label{saut:eq:ell}
 \ell(g)=[g]_0,
\end{equation}
its real coefficient at exponent zero. This is an additive map defined on
all of $\No$, and $\ell(r)=r$ for every real $r$. It is not a standard-part
map on all surreals: for an infinite $g$, coefficient extraction still makes
sense even though standard part does not.

\subsection{Complex phases move the real axis}

For $\theta\in\R$ use the same $\ell$ to define
\begin{equation}\label{saut:eq:phase}
 P_\theta\left(\sum_g z_g t^g\right)
 =\sum_g z_g\exp(i\theta\ell(g))t^g.
\end{equation}
Only the ordinary complex exponential is used in the coefficient.

\begin{theorem}[A faithful phase action]\label{saut:thm:phase}
The maps $P_\theta$ form a faithful action of the additive group $\R$ by
strongly $\C$-linear, exactly value-preserving automorphisms of $\K$.
For $\theta\ne0$, $P_\theta$ does not preserve $\No$. Moreover,
\begin{equation}\label{saut:eq:phaseconj}
 P_\theta P_\phi=P_{\theta+\phi},\qquad
 cP_\theta c=P_{-\theta}.
\end{equation}
In particular,
\[
 P_{\pi/2}(t)=it.
\]
\end{theorem}
\begin{proof}
The character identity follows from additivity of $\ell$, so the
monomial theorem gives all field and summability assertions. If
$P_\theta=\id$, then $e^{i\theta r}=1$ for every real $r$, because
$\ell(r)=r$. For $\theta\ne0$, taking $r=\pi/\theta$ contradicts this.
Thus the action is faithful. Its parameter is not an angle modulo $2\pi$:
all real exponents occur, not just integer exponents.

If $\theta\ne0$, the choice $r=\pi/(2\theta)$ yields
$P_\theta(t^r)=it^r$, which is not real. Finally, conjugation reverses every
coefficient phase, giving \eqref{saut:eq:phaseconj}.
\end{proof}

This provides the second strict inclusion in \eqref{saut:eq:three} without a
choice of transcendence basis or a nonconstructive extension. It also proves
that none of the following conditions by itself recovers the real axis:
fixing ordinary $\C$, strong $\C$-linearity, exact valuation preservation,
or even all three together.

Valuation preservation is much weaker than exact modulus preservation.
For example,
\[
 |1+t|=1+t,
 \qquad |P_{\pi/2}(1+t)|=|1+it|=\sqrt{1+t^2}\ne1+t.
\]
Likewise, $P_{\pi/2}$ does not preserve the distinguished unit circle,
by \cref{saut:thm:circle}. Algebraic multiplication survives while the chosen
Euclidean geometry changes.

\subsection{A continuum of explicit conjugations}

Set
\begin{equation}\label{saut:eq:twistedconj}
 j_\theta=P_\theta cP_\theta^{-1}=P_{2\theta}c,
 \qquad F_\theta=P_\theta(\No).
\end{equation}
Then $j_\theta^2=\id$, $\Fix(j_\theta)=F_\theta$, and
$\K=F_\theta(i)$. Each $F_\theta$ is a real closed subfield isomorphic to
$\No$, but the actual embedded real axes are pairwise distinct.
Indeed, equality $F_\theta=F_\phi$ would make $P_{\theta-\phi}$ preserve
$\No$, forcing $\theta=\phi$. The fixed-field condition is especially
concrete:
\begin{equation}\label{saut:eq:twistedcoeff}
 z\in F_\theta\quad\Longleftrightarrow\quad
 z_g e^{-i\theta\ell(g)}\in\R\quad\text{for all }g.
\end{equation}
These are different \emph{placements} of an isomorphic real form. We will
also construct a real form not isomorphic to $\No$ in
Section~\ref{saut:sec:realforms}.

\section{Leading-term kernels and infinitesimal motion}\label{saut:sec:kernels}

Define
\begin{equation}\label{saut:eq:U}
 \U=\{\alpha\in\Aut(\K):
       v(\alpha z-z)>v(z)\text{ for all }z\ne0\}.
\end{equation}
Here $v(0)=\infty$, so fixed inputs satisfy the inequality. Equivalently,
$\alpha$ fixes the complete leading term of every nonzero series.
Such maps are often called $1$-automorphisms. They can still change very
substantial lower-order data. Let $\U_{\R}$ denote the corresponding
leading-term-fixing subgroup of $\Aut(\No)$.

\subsection{A real automorphism need not fix ordinary reals}

An order-preserving automorphism $\sigma$ of $\No$ fixes every rational.
For $r\in\R$, rational comparisons imply only
\[
 \sigma(r)-r\in\mm\cap\No,
\]
not, for a general $r$, $\sigma(r)=r$. For an \emph{irrational} real $r$,
every infinitesimal perturbation has the same strict rational comparisons:
each nonzero ordinary distance $|r-q|$, $q\in\Q$, dominates every
infinitesimal. Rational $r$ are fixed already, and real algebraic $r$ are
fixed by their polynomial and order data (\cref{saut:lem:order}).
The following construction makes the distinction explicit.

\begin{theorem}[Taylor motion of coefficient constants]\label{saut:thm:coeffflow}
Let $d:\C\to\C$ be a field derivation and let $\delta\in\Gamma$ be
positive. Then
\begin{equation}\label{saut:eq:coeffflow}
 \Phi_d\left(\sum_g z_gt^g\right)
 =\sum_{g,\,n\ge0}\frac{d^n(z_g)}{n!}t^{g+n\delta}
\end{equation}
defines a strong automorphism belonging to $\U$, with inverse $\Phi_{-d}$.
It fixes all monomials $t^g$. If $d(\R)\subseteq\R$ and $d(i)=0$, it
preserves $\No$ and commutes with conjugation.
There exist choices for which a real transcendental $b$ satisfies
\begin{equation}\label{saut:eq:realmove}
 \Phi_d(b)=b+t^\delta.
\end{equation}
\end{theorem}
\begin{proof}
For an input with support $S$, all output exponents lie in $S+\N\delta$.
Lemma~\ref{saut:lem:support} proves well-ordering and finite incidence at every
exponent. The same argument works for a summable family of inputs and for
all double sums needed in a composition.
The iterated Leibniz rule
\[
 d^n(ab)=\sum_{k=0}^n\binom nk d^k(a)d^{n-k}(b)
\]
proves multiplicativity after division by $n!$. Additivity is immediate.
In $\Phi_d\Phi_{-d}$, the coefficient multiplying the $k$th derivative at total
shift $k\delta$ is
\[
 \sum_{n=0}^k\frac{(-1)^n}{n!(k-n)!},
\]
which is $1$ for $k=0$ and $0$ otherwise. This proves the inverse formula.
Every term with $n>0$ has exponent strictly larger than the original
minimum; the $n=0$ term retains its original leading coefficient. Hence
$\Phi_d\in\U$. Since $d(1)=0$, each monomial is fixed.
For any $z\ne0$, the correction in fact satisfies
$v(\Phi_d(z)-z)\ge v(z)+\delta$, with $v(0)=\infty$.
Also $d(i)=0$ is automatic from $d(i^2+1)=2i\,d(i)=0$; preservation of
the real coefficients is the extra condition in the real-axis assertion.

To obtain \eqref{saut:eq:realmove}, choose a transcendence basis $B$ of $\R/\Q$
and an element $b\in B$. Define a derivation on $\Q(B)$ by $d(b)=1$ and
$d(b')=0$ for $b'\ne b$. In characteristic zero it extends uniquely through
an algebraic extension: for a minimal polynomial $p(X)$ and a root $a$,
\[
 d(a)=-\frac{p^d(a)}{p'(a)},
\]
where $p^d$ means differentiate the coefficients. Extend first to $\R$,
then to $\C$ by $d(i)=0$. Now $d^2(b)=0$, giving \eqref{saut:eq:realmove}.
All real coefficients stay real, and the construction commutes with $c$.
\end{proof}

Thus even a strongly additive, leading-term-fixing surreal automorphism
can move an ordinary real constant. It is not strongly $\R$-linear.
Confusing strong additivity with coefficient linearity would erase this
entire type of automorphism.

The construction uses choice only for a derivation on the \emph{set} $\R$;
no class-sized basis is selected. Its action on a given normal form is then
an explicit support-controlled formula. It is generally noncomputable from
ordinary numerical names of the real coefficients, as should be expected
from its use of a transcendence basis.

\subsection{Strongly complex-linear flows}

There are also abundant nontrivial members of $\U$ fixing all constants.
Let $\lambda:\Gamma\to\C$ be additive, and fix $\delta>0$. Define
\begin{equation}\label{saut:eq:shiftderiv}
 \mathcal D\left(\sum_g z_gt^g\right)=\sum_g z_g\lambda(g)t^{g+\delta}.
\end{equation}
The positive shift is important: $g-\delta$ would generally lower valuation,
not raise it.

\begin{theorem}[A fixed-shift flow]\label{saut:thm:shiftflow}
The map $\mathcal D$ is a strongly $\C$-linear derivation, and for every $s\in\C$,
\begin{equation}\label{saut:eq:shiftflow}
 T_s=\sum_{n\ge0}\frac{s^n\mathcal D^n}{n!}
\end{equation}
is a strongly $\C$-linear automorphism in $\U$. The group law is
$T_sT_r=T_{s+r}$. If $\lambda$ is real-valued and $s\in\R$, then $T_s$
commutes with $c$ and preserves $\No$.
\end{theorem}
\begin{proof}
Additivity of $\lambda$ proves the Leibniz rule on monomials, and finite
convolution proves it on all series. Iterating gives
\[
 \mathcal D^n(t^g)=
 \left(\prod_{k=0}^{n-1}\lambda(g+k\delta)\right)t^{g+n\delta}.
\]
All supports in \eqref{saut:eq:shiftflow} lie in $S+\N\delta$. The support lemma
justifies summation and the finite regroupings at each coefficient.
The iterated Leibniz identity proves multiplicativity, and the binomial
formula gives $T_sT_r=T_{s+r}$. Thus $T_{-s}$ is the inverse. Nonconstant
terms of the operator exponential strictly increase every input's valuation,
so the leading term is unchanged. Constants are killed by $\mathcal D$ because
$\lambda(0)=0$. The real-coefficient assertion is immediate.
\end{proof}

For a particularly transparent example choose $\delta=1$ and $\lambda=\ell$
from \eqref{saut:eq:ell}. Then
\[
 \mathcal D(t)=t^2,\qquad \mathcal D^n(t)=n!t^{n+1},\qquad
 T_s(t)=\frac{t}{1-st}.
\]
Consequently $T_1$ is a nonidentity strongly $\C$-linear, real-axis-preserving
$1$-automorphism. These explicit flows illustrate the automorphism--derivation
connection without requiring a blanket assertion that every formally
written operator exponential is summable.

\subsection{What monomial prescriptions do not automatically give}

For a strongly $\C$-linear element $\alpha\in\U$, the assignments
\[
 u_g=t^{-g}\alpha(t^g)\in1+\mm
\]
form a group homomorphism $\Gamma\to1+\mm$. Formally one would like to write
\[
 \alpha\left(\sum_g z_gt^g\right)=\sum_g z_gt^g u_g.
\]
This is valid for an existing strong automorphism. Conversely, an arbitrary
homomorphism $g\mapsto u_g$ does not by itself certify that this formula is
summable for \emph{every} well-ordered support or that it is onto.
One must check the transformed support family and construct a compatible
inverse. The fixed-shift proofs above supply exactly those missing checks.
This is a genuine obstacle in turning a group decomposition into a complete
explicit parametrization of its largest kernel.

\section{Decomposing the valued automorphism group}\label{saut:sec:decomposition}

Define
\[
 G_v=\{\alpha\in\Aut(\K):\alpha(\OO)=\OO\}.
\]
Preserving the valuation here means preserving its comparison structure,
not necessarily each value. Such an $\alpha$ induces
\[
 \rho_\alpha\in\Aut(\C),\qquad
 \tau_\alpha\in\Gh,
 \qquad
 v(\alpha z)=\tau_\alpha(vz).
\]
The residue action is $\rho_\alpha(\st z)=\st(\alpha z)$ for finite $z$.
This does not assert $\alpha(\C)=\C$: the coefficient-motion examples have
identity residue action while sending constants out of the constant field.

\begin{theorem}[Four layers for surcomplex valued automorphisms]\label{saut:thm:decomp}
With $\U$ as in \eqref{saut:eq:U}, every $\alpha\in G_v$ factors uniquely as
\begin{equation}\label{saut:eq:factor}
 \alpha=u\,D_\chi\,M_{\rho,\tau},
\end{equation}
where $u\in\U$, $\chi\in\Hom(\Gamma,\C^\times)$, and
\[
 D_\chi\left(\sum z_gt^g\right)=\sum z_g\chi(g)t^g,
 \qquad
 M_{\rho,\tau}\left(\sum z_gt^g\right)=\sum\rho(z_g)t^{\tau(g)}.
\]
(In the notation of \eqref{saut:eq:monomial}, $D_\chi=M_{\id,\id,\chi}$ and
$M_{\rho,\tau}=M_{\rho,\tau,1}$; the maps $D_\lambda$ of
Section~\ref{saut:sec:explicit} are the $D_\chi$ with $\chi=\exp\circ\lambda$.)
In group notation,
\begin{equation}\label{saut:eq:fourfactor}
 G_v\simeq
 \bigl(\U\rtimes\Hom(\Gamma,\C^\times)\bigr)
 \rtimes\bigl(\Aut(\C)\times\Gh\bigr).
\end{equation}
The outer action on the character factor is
\begin{equation}\label{saut:eq:charouter}
 ((\rho,\tau)\cdot\chi)(g)=\rho\bigl(\chi(\tau^{-1}g)\bigr).
\end{equation}
The other actions are conjugation by the displayed lifts.
\end{theorem}
\begin{proof}
Stabilizing $\OO$ also stabilizes its units and its maximal ideal, so it
induces the residue automorphism $\rho_\alpha$. Define directly
$\tau_\alpha(g)=v(\alpha(t^g))$. If $v(z)=g$, then $z/t^g$ is a unit,
so $v(\alpha z)=\tau_\alpha(g)$. This proves independence from the chosen
representative. Multiplication proves additivity; membership of quotients
in $\OO$ proves order preservation; the same construction for
$\alpha^{-1}$ supplies the inverse. Thus no quotient whose elements would
be proper classes needs to be formed.
The maps $M_{\rho,\tau}$ give a section of the induced-action homomorphism,
by the monomial construction. Let $I$ be its kernel and put
$\beta=\alpha M_{\rho,\tau}^{-1}\in I$.

For $\beta\in I$, define
\[
 \chi_\beta(g)=\lc(\beta(t^g)).
\]
Since $\beta$ fixes values, multiplicativity of leading coefficients gives
$\chi_\beta(g+h)=\chi_\beta(g)\chi_\beta(h)$. For an arbitrary nonzero $z$,
write $z=a t^g(1+\varepsilon)$ with $a=\lc z$ and $v\varepsilon>0$.
Identity residue action gives $\lc(\beta a)=a$, and
$\beta(1+\varepsilon)$ still has leading coefficient $1$. Hence
\begin{equation}\label{saut:eq:internalLC}
 \lc(\beta z)=\lc(z)\chi_\beta(vz).
\end{equation}
It follows that $\beta\mapsto\chi_\beta$ is a group homomorphism on $I$,
with kernel exactly $\U$. The maps $D_\chi$ split it. Thus
$u=\beta D_{\chi_\beta}^{-1}$ gives \eqref{saut:eq:factor}, and all parameters
are forced by the preceding operations. More explicitly, the character in
that factorization is
$\chi(g)=\lc(\alpha(t^{\tau^{-1}(g)}))$; it is indexed after undoing the
value action. The subgroup $\U$ is normal in $G_v$, since conjugation by
a valued automorphism transports its strict valuation inequality through
the ordered map $\tau$. Evaluating
$M_{\rho,\tau}D_\chi M_{\rho,\tau}^{-1}$ on a monomial gives
\eqref{saut:eq:charouter}.
\end{proof}

This direct argument is the full-Hahn instance of the valued-field structure
studied in \cite{KS}. It also clarifies its limits: \eqref{saut:eq:fourfactor}
is not a decomposition of \emph{all} surcomplex field automorphisms. We will
exhibit automorphisms outside $G_v$ in Section~\ref{saut:sec:bare}.

The same proof works with the strong subgroup, replacing $\U$ by its strong
part. All explicit sections are strong, so strength of an automorphism
is concentrated in the residual kernel factor. The factor $\U$ itself is
not replaced by a conveniently finite list of coefficients. A strong field
map is determined by its action on \emph{both} ordinary coefficients and
monomials, using
$\alpha(\sum z_gt^g)=\sum\alpha(z_g)\alpha(t^g)$.
Monomial images alone suffice when it also fixes $\C$ pointwise.
Strongness alone does not remove the coefficient data: the nonidentity
Taylor motions \cref{saut:thm:coeffflow} fix every monomial. Without
strongness, the displayed reconstruction is not licensed:
Section~\ref{saut:sec:topology} gives an automorphism that fixes every
summand of a Hahn family but moves its sum.

\subsection{The real reduction}

Every field automorphism of $\No$ preserves order, hence preserves the
natural valuation ring. Its residue action on $\R$ is an ordered-field
automorphism and therefore identity. Applying the same proof over $\R$
gives
\begin{equation}\label{saut:eq:realdecomp}
 \Aut(\No)\simeq
 \bigl(\U_{\R}\rtimes\Hom(\Gamma,\R_{>0})\bigr)\rtimes\Gh.
\end{equation}
Every character into $\R^\times$ is positive because $\Gamma$ is divisible.
Applying the ordinary logarithm pointwise identifies its group with
$\Hom(\Gamma,\R)$ under addition.

Combining \eqref{saut:eq:realdecomp} with the real-axis theorem gives a structural
description of $G_{\No}$:
\begin{equation}\label{saut:eq:realaxisdecomp}
 G_{\No}\simeq
 \left[\bigl(\U_{\R}\rtimes\Hom(\Gamma,\R_{>0})\bigr)\rtimes\Gh\right]
 \times C_2.
\end{equation}
In particular, ``residue action is identity'' must not be shortened to
``all embedded real constants are fixed.'' Theorem~\ref{saut:thm:coeffflow}
is an explicit counterexample to that shortcut.

\subsection{One level deeper: the ordered additive value group}

Because $\Gamma=(\No,+,<)$ is itself the ordered additive group of real
Hahn series, one can further analyze $\Gh$. Let $\U_{\mathrm{add}}$ consist
of its additive automorphisms fixing each leading term. For an order
bijection $\zeta:\Gamma\to\Gamma$ and positive real scalars
$(a_g)_{g\in\Gamma}$, the formula
\begin{equation}\label{saut:eq:skeleton}
 A_{\zeta,a}\left(\sum_g r_gt^g\right)
 =\sum_g a_g r_g t^{\zeta(g)}
\end{equation}
defines an ordered additive automorphism. Here $\zeta$ need not be additive:
this is an additive-group construction, not a field construction.

An ordered additive automorphism $\tau$ permutes the Archimedean classes
in order. Its parameters can be read directly as
$\zeta(g)=v(\tau(t^g))$ and $a_g=\lc(\tau(t^g))>0$.
For fixed $g$, comparison of $rt^g$ with rational multiples of $t^g$
shows that the leading coefficient of $\tau(rt^g)$ is $a_gr$ for every
$r\in\R^\times$, at value $\zeta(g)$. A tail of value greater than $g$
is infinitesimal relative to $t^g$ in the ordered additive group, and its
image has value greater than $\zeta(g)$. Thus the leading term of the
image of any input with leading term $rt^g$ is
$a_grt^{\zeta(g)}$. This uses only order and additivity, not strong
summation or a quotient whose elements would be proper classes.
Removing \eqref{saut:eq:skeleton} leaves exactly
$\U_{\mathrm{add}}$, and the parameters are unique. This is the additional
kernel--skeleton factorization, parallel to \cite{KSG}.

It is important to retain the reindexing action in its group law. Direct
calculation gives
\begin{equation}\label{saut:eq:skeletonlaw}
 A_{\zeta,a} A_{\eta,b}=A_{\zeta\eta,d},
 \qquad d_g=a_{\eta(g)}b_g.
\end{equation}
Thus these parameters have a semidirect-product law, not componentwise
multiplication of an unrelated permutation and weight family. The ordinary
logarithm turns positive weights into additive real weights but does not
remove their reindexing. Together with \eqref{saut:eq:realdecomp}, this identifies
five structural layers on the real side: field leading-term motion,
characters, additive leading-term motion, exponent-order permutations,
and Archimedean coefficient weights.

\section{Fine topology does not mean strong summation}\label{saut:sec:topology}

The fine topology on $\K$ is generated by the balls
\[
 B(a,r)=\{z:|z-a|<r\},\qquad r\in\No_{>0}.
\]
By \eqref{saut:eq:normvalue}, it is the same topology as the valuation topology
with neighborhoods
\[
 V_\gamma(a)=\{z:v(z-a)>\gamma\},\qquad\gamma\in\Gamma.
\]
For example, $V_{v(r)}(a)\subseteq B(a,r)$, while a ball of radius
$t^{\gamma+1}$ is contained in $V_\gamma(a)$. These inclusions suffice to
identify the neighborhood systems.

\begin{proposition}\label{saut:prop:continuous}
Every element of $G_v$ is a fine homeomorphism and a uniform equivalence.
In particular, every real-axis-preserving surcomplex field automorphism is
fine-continuous.
\end{proposition}
\begin{proof}
If $\alpha$ induces $\tau$ on $\Gamma$, then
\[
 v(\alpha z-\alpha a)=\tau(v(z-a)).
\]
The order bijection $\tau$ takes the neighborhood basis to itself: for a
target threshold $\eta$, use the input threshold $\tau^{-1}(\eta)$.
This threshold is independent of the base point. Applying
the same equation to pairs of points proves uniform continuity of both
$\alpha$ and its inverse. The real-axis subgroup lies in $G_v$ by
\cref{saut:lem:order} and \eqref{saut:eq:equivariantnorm}.
\end{proof}

This is an inclusion in the continuous automorphism group, not a
classification of that whole group. In particular, we do not infer a
converse merely from continuity. The chosen valuation ring contains
additional information beyond the neighborhood calculation used here.

The phase twists are therefore continuous even though they do not preserve
the real axis. Lifting any ordinary field automorphism $\rho$ of $\C$
coefficientwise gives further examples in $G_v$. There is no contradiction
with the familiar classification of continuous automorphisms of ordinary
$\C$: the topology induced on $\C\subset\K$ is \emph{discrete}, not its
usual Archimedean topology. A ball of infinitesimal radius around a constant
contains no other ordinary complex constant.

\subsection{Why small nets cannot test the full topology}

\begin{proposition}\label{saut:prop:discrete}
Every set-sized subset of the full surcomplex field is closed and discrete.
Every convergent set-indexed net is eventually equal to its limit.
Nevertheless, the full field has no isolated points.
\end{proposition}
\begin{proof}
For a set $S\subseteq\K$ and a point $a$, the values
$\{v(s-a):s\in S,\ s\ne a\}$ form a set of surreal numbers. Every set of
surreals has an upper bound, so choose $\gamma$ strictly above all of them.
Then $V_\gamma(a)$ misses $S\setminus\{a\}$. This proves both closedness
and discreteness. Apply the same argument to the range of a set-indexed
net together with its limit to get eventual equality. Finally,
$a+t^{\gamma+1}\ne a$ belongs to $V_\gamma(a)$, so no point of the full
field is isolated.
\end{proof}

For example, $t^n$ does not tend to zero in the full surreal fine topology:
none of these terms lies in $V_\omega(0)$. In the fixed field
$\C((t^\Q))$, by contrast, the integers are cofinal in its value group
and $t^n\to0$ in its intrinsic valuation topology.
A set-sized Hahn workspace inside $\K$ acquires a discrete subspace
topology from the full field, even when its intrinsic valuation topology
is nondiscrete. The choice of topology is part of the theorem, not merely
a presentational convention.

\subsection{A continuous automorphism that fails to preserve a Hahn sum}

We give an explicit support-correct witness. We will use the standard
Newton--Puiseux fact that
\[
 \mathcal P=\bigcup_{m\ge1}\R((t^{1/m}))
\]
is real closed. Indeed,
$\mathcal P(i)=\bigcup_{m\ge1}\C((t^{1/m}))$ is algebraically closed by
Newton--Puiseux \cite{ParanVo}; an ordered field with algebraically closed
quadratic extension by $i$ is real closed. This set-sized Puiseux field is
also discussed among the special cases in \cite[Section~5.5]{KS}. Define
\begin{equation}\label{saut:eq:nonstronga}
 a=\sum_{n\ge2}t^{1-1/n},\qquad \varepsilon=t^\omega.
\end{equation}
The support $\{1/2,2/3,3/4,\ldots\}$ is a well-ordered increasing sequence.
Its denominators are unbounded, so $a\notin \mathcal P$. Since $\mathcal P$ is real closed,
$a$ is transcendental over $\mathcal P$.

For each $p\in \mathcal P$, the nonzero difference $a-p$ has rational support and
hence rational valuation. Thus $\varepsilon$ is smaller than $|a-p|$.
It follows that $a$ and $a+\varepsilon$ have the same cut over $\mathcal P$.
An ordered isomorphism $\mathcal P(a)\to \mathcal P(a+\varepsilon)$ fixing $\mathcal P$ and moving $a$
therefore exists; real closedness of $\mathcal P$ ensures that signs of rational
functions are determined by this cut. Extend it through real closures and
then to an automorphism $\sigma$ of $\No$ by the ordered extension theorem
proved in Section~\ref{saut:sec:bare}.

Every summand in \eqref{saut:eq:nonstronga} belongs to $\mathcal P$ and is fixed. But
\[
 \sigma(a)=a+\varepsilon\ne a
     =\sum_{n\ge2}\sigma(t^{1-1/n}).
\]
The lift of $\sigma$ to $\K$ fixes ordinary $\C$, commutes with conjugation,
and is fine-continuous, but is not strong.
We have not claimed that this particular extension fixes every value or
every leading term; those stronger assertions are unnecessary here.

There is also a finite-observation obstruction to approximation by strong
maps. Any real surreal automorphism fixing $t$ fixes all its positive
rational powers, by uniqueness of positive roots. If it is strong, it fixes
$a$. Hence the neighborhood specified by
\[
 \psi(t)=t,\qquad \psi(a)=a+\varepsilon
\]
in the pointwise-equality topology on $\Aut(\No)$ contains $\sigma$ and
contains no strong automorphism. This proves non-density in that particular
automorphism-group topology, not in every possible topology on the group.

The same finite observations exclude strong maps in the full $\Aut(\K)$.
If $\psi(t)=t$ and $q_n=1-1/n$, then
$\psi(t^{q_n})^n=t^{n-1}$, so
$\psi(t^{q_n})=\zeta_n t^{q_n}$ with $\zeta_n$ an ordinary complex
$n$th root of unity. Strongness would force
$\psi(a)=\sum_{n\ge2}\zeta_n t^{q_n}$, whose support is contained in
$\Q$ and cannot contain the exponent $\omega$ of $a+\varepsilon$.
The lifted witness satisfies both observations, so the strong subgroup
is not dense there either. This argument requires neither preservation of
the real axis nor fixation of all ordinary complex constants by $\psi$.

\section{Differentiability: a surprising zero-derivative phenomenon}\label{saut:sec:derivative}

For a function $f:\K\to\K$, its fine derivative at $z$ is the
punctured-neighborhood limit
\[
 f'(z)=\lim_{h\to0,\ h\ne0}\frac{f(z+h)-f(z)}h,
\]
when this limit exists in the full fine topology. The limit quantifies over
all small surcomplex increments; it is not tested only on an ordinary
sequence. This is the same difference-quotient convention encoded in the
repository's generic \texttt{FineHasDerivAt} interface \cite{RepoFine}.

\begin{theorem}[Derivative dichotomy for field automorphisms]\label{saut:thm:derivative}
If a field automorphism $\alpha$ of $\K$ has a fine derivative at one point,
it has the same derivative at every point. If that derivative is nonzero,
then $\alpha=\id$ and the derivative is $1$.
Nonidentity differentiable automorphisms are nevertheless possible: they
must have derivative $0$ everywhere.
\end{theorem}
\begin{proof}
Additivity makes every difference quotient equal to
$q(h)=\alpha(h)/h$, independently of the base point. For a fixed $b\ne0$,
\[
 q(bh)=\frac{\alpha(b)}bq(h).
\]
If $q(h)\to L$, multiplication by $b$ is a homeomorphism of the punctured
neighborhood of zero, so the left-hand side also tends to $L$.
Hausdorff uniqueness gives $L=\alpha(b)L/b$. When $L\ne0$, cancellation
forces $\alpha(b)=b$ for every nonzero $b$. Thus $\alpha=\id$, whose
quotient is constantly $1$.
\end{proof}

\begin{theorem}[Exponent dilations have derivative zero]\label{saut:thm:zeroderiv}
For the maps $S_a$ of \eqref{saut:eq:dilation},
\[
 S_a'\equiv0\quad(a>1),\qquad S_1'\equiv1.
\]
For $0<a<1$, $S_a$ has no $\K$-valued fine derivative at any point.
In particular, $S_2$ is a nonidentity field automorphism and a fine
homeomorphism with derivative zero everywhere.
\end{theorem}
\begin{proof}
For every nonzero $h$,
\begin{equation}\label{saut:eq:dilationquot}
 v\left(\frac{S_a(h)}h\right)=(a-1)v(h).
\end{equation}
If $a>1$, given any $\eta\in\Gamma$, taking
$v(h)>\eta/(a-1)$ forces the quotient into $V_\eta(0)$.
Thus the quotient tends to zero. This argument works for positive
infinitesimal $a-1$ as well as for ordinary positive $a-1$.
If $0<a<1$, \cref{saut:thm:derivative} would force any derivative to be
zero, since $S_a(t)=t^a\ne t$. But for every proposed input threshold
$\gamma_0$, choose $\gamma>\max\{\gamma_0,0\}$ and put $h=t^\gamma$.
Then $h\in V_{\gamma_0}(0)$, while the quotient has value
$(a-1)\gamma<0$ and lies outside $V_0(0)$. This rules out a zero limit
using the full neighbourhood criterion, without a set-indexed cofinal net.
\end{proof}

Conjugating the input does not change its valuation. Thus $S_a c$ also
has zero derivative for $a>1$, even though $c$ by itself is not
surcomplex-differentiable. For $c$, real and imaginary increments give
quotients $1$ and $-1$ at arbitrarily small scales.

\begin{corollary}\label{saut:cor:bidifferentiable}
If a field automorphism and its inverse are both fine-differentiable, the
automorphism is identity.
\end{corollary}
\begin{proof}
Differentiability implies continuity by writing an increment of the
function as its difference quotient times the input increment. If both
maps were nonidentity and differentiable, both derivatives would be zero.
Put $k=\alpha(h)$; continuity implies $k\to0$ as $h\to0$. The product
\[
 \frac{\alpha(h)}h\,
 \frac{\alpha^{-1}(k)}k=1
\]
would then tend to zero, a contradiction. Injectivity ensures $k\ne0$
for $h\ne0$, so the inverse's punctured limit applies. The identity satisfies both
conditions.
\end{proof}

These results do not assert that $S_2$ is representable by a locally
convergent power series with fixed coefficients. They concern the fine
difference-quotient definition. In this setting, ``everywhere
differentiable'' cannot be used as a synonym for ordinary complex analytic
rigidity, and a zero derivative does not imply local constancy: $S_2$ is
injective and the full field has no isolated points. The
Berarducci--Mantova scalar derivation discussed later is a different
operation again.

\section{The unrestricted field: homogeneity and fixed fields}\label{saut:sec:bare}

The preceding explicit formulas already disprove real-axis rigidity.
The full pure-field picture is even less rigid. To discuss it without
misusing a set-sized extension theorem on a proper class, we give the
extension argument directly.

\begin{theorem}[Class back-and-forth]\label{saut:thm:acfhomogeneity}
Assume global choice. Let $K_1,K_2$ be proper-class algebraically closed
fields of the same characteristic. Every isomorphism between set-sized
subfields of $K_1$ and $K_2$ extends to an isomorphism $K_1\to K_2$.
In particular, every isomorphism between set-sized subfields of $\K$
extends to a field automorphism of $\K$.
\end{theorem}
\begin{proof}
Fix a set-like global well-order and the resulting class enumerations of
$K_1$ and $K_2$, indexed by $\On$.
Recursively maintain an isomorphism $f_\alpha:A_\alpha\to B_\alpha$
between \emph{set-sized} subfields. At a forth step, take the next element
$x$ of $K_1$. If $x$ is algebraic over $A_\alpha$, send it to a root in
$K_2$ of its transported irreducible polynomial, extending the field
isomorphism. If $x$ is transcendental, choose a transcendental element over
$B_\alpha$. This is possible because the elements of $K_2$ algebraic over
a set-sized subfield form a set: there are set many polynomials and each
has finitely many roots. A proper-class field cannot be exhausted by it.

At each successor stage perform both a forth step and a back step, the
latter for the next element of $K_2$ using the inverse partial map; skip an
element already in the relevant partial field. Always select the least
admissible witness in the fixed global well-order. At a limit ordinal
take unions. There are only set many
earlier stages at any given ordinal, so the partial domains and ranges
remain sets. The class union of all partial maps is a field isomorphism;
forth steps exhaust $K_1$ and back steps exhaust $K_2$.
Each partial map is a set. The deterministic rule has a unique recursion
of every set ordinal length, with the ambient fields and well-order as
class parameters; these compatible recursions define the class union.
The candidate witnesses can form a proper class, but its least member is
selected by the global well-order. No recursion whose stage value is itself
a proper class, or class-sized transcendence basis, is needed.
\end{proof}

\begin{theorem}[Ordered surreal extension]\label{saut:thm:orderedhomogeneity}
Under the same foundational convention, every ordered-field isomorphism
between set-sized subfields of $\No$ extends to an automorphism of $\No$.
\end{theorem}
\begin{proof}
First extend the map to real closures inside $\No$; uniqueness of ordered
real closures supplies the extension. Run a back-and-forth as above, now
keeping the partial fields real closed. An element $x$ outside a partial
real closed field $A$ is transcendental over it. Its comparisons with $A$
form a set-sized cut $(\mathcal L,\mathcal R)$ with $\mathcal L\cup\mathcal R=A$.
Transport this cut to the current range $B$. The surreal cut property
supplies an element $y$ strictly between every element of its left and
right sides. It cannot belong to $B$, since every element of $B$ occurs
on one of those sides. Thus $y$ is transcendental over $B$.

The substitution $x\mapsto y$ gives an ordered isomorphism of rational
function fields. To see preservation of signs, factor polynomials over
the real closed coefficient field into linear factors and quadratic
factors without real roots; their signs are determined by the cut and
leading coefficient. Extend again to real closures. At limits, unions of
increasing real closed fields are real closed, since each polynomial has
its finitely many coefficients in one stage. Forth and back exhaust $\No$.
\end{proof}

These proofs explain why field homogeneity is so strong and why it should
not be confused with simplicity preservation. The selected separator $y$
need not have the same birthday or sign-sequence complexity as $x$.

\subsection{Orbits over ordinary complex constants}

Every element of $\K\setminus\C$ is transcendental over $\C$, since $\C$
is algebraically closed. Therefore, for $x,y\notin\C$, the substitution
$x\mapsto y$ is an isomorphism $\C(x)\to\C(y)$ and extends by
\cref{saut:thm:acfhomogeneity}. We obtain
\begin{equation}\label{saut:eq:transitive}
 \Aut(\K/\C)\text{ is transitive on }\K\setminus\C.
\end{equation}
The analogous statement for finite algebraically independent tuples follows
by the same rational-function argument. Thus a pure-field automorphism
fixing every ordinary complex number can send an infinite surreal to an
infinitesimal, or to an element off the real axis.

\begin{proposition}\label{saut:prop:nondefinability}
For every set $A\subseteq\K$, there are automorphisms fixing $A\cup\C$
pointwise which do not preserve $\No$. There are also such automorphisms
which do not preserve $\OO$.
Consequently, neither the distinguished real axis nor the natural valuation
ring is invariant under all pure-field automorphisms over any prescribed
set of parameters.
\end{proposition}
\begin{proof}
Let $E$ be the algebraic closure inside $\K$ of the field generated by
$A\cup\C$. It is a set. Choose a positive infinite $x\in\No\setminus E$;
for instance choose $x$ larger than the moduli of all elements of $E$.
Then $x$ is transcendental over $E$. Both maps
\[
 E(x)\longrightarrow E(ix),\quad x\mapsto ix,
 \qquad
 E(x)\longrightarrow E(x^{-1}),\quad x\mapsto x^{-1}
\]
are field isomorphisms fixing $E$. Extend each to $\K$. The first moves a
real element off the real axis. The second sends an element outside $\OO$
into its maximal ideal, so cannot stabilize $\OO$.
\end{proof}

In particular, these structures are not first-order definable in the pure
field with parameters: definable sets are invariant under automorphisms
fixing their parameters. This conclusion is stronger than merely saying
that no preferred coordinate system has yet been chosen.
Combining the proposition with the phase example gives the proper chain
\begin{equation}\label{saut:eq:strictchain}
 G_{\No}\ \subsetneq\ G_v\ \subsetneq\ \Aut(\K).
\end{equation}
The first strictness has an explicit strong, $\C$-linear, value-fixing witness;
the second has a class back-and-forth witness.

\subsection{Fixed fields of the main collections}

Let $\Q^{\rc}\subset\R$ be the field of real algebraic numbers and let
$\overline\Q=\Q^{\rc}(i)\subset\C$.

\begin{proposition}\label{saut:prop:fixedfields}
In the class-map interpretation,
\begin{align*}
 \Fix(\Aut(\K))&=\Q, &
 \Fix(\Aut(\K/\C))&=\C,\\
 \Fix(\Aut(\No))&=\Q^{\rc}, &
 \Fix(G_{\No})&=\Q^{\rc}.
\end{align*}
The fixed field of the $i$-fixing lifts of all surreal automorphisms is
$\overline\Q$.
\end{proposition}
\begin{proof}
A nonrational algebraic element has a distinct conjugate in characteristic
zero, and its conjugacy isomorphism extends by
\cref{saut:thm:acfhomogeneity}. A transcendental element can similarly be moved
by a rational-function isomorphism. This gives the first assertion.
Transitivity \eqref{saut:eq:transitive} gives the second.

Every surreal automorphism fixes $\Q^{\rc}$ by order. For
$x\notin\Q^{\rc}$, $x$ is transcendental over this real closed field.
Choose a positive $\varepsilon$ smaller than all nonzero distances
$|x-a|$ for $a\in\Q^{\rc}$. This is possible for a set of distances by
the surreal cut property. Then $x$ and $x+\varepsilon$ have the same cut,
so the ordered extension theorem moves $x$. This proves the third
assertion. For the last two assertions apply the coordinate formula
\eqref{saut:eq:lift}; allowing $c$ kills the imaginary fixed coordinate,
whereas forbidding $c$ leaves both coordinates in $\Q^{\rc}$.
\end{proof}

\subsection{Why a transcendence-basis slogan is not a classification}

For a set-sized algebraically closed field, specifying a transcendence
basis is useful for producing automorphisms: first choose compatible images
of transcendental generators and then extend through algebraic elements.
But an automorphism need not permute one fixed basis, and its action on
algebraic elements is additional data. Even there, ``the automorphism group
is the symmetric group on a transcendence basis'' is false.

For the full surcomplex field, the corresponding basis would be a proper
class and requires additional foundational care. The set-stage
back-and-forth proof is sufficient for all existential statements here.
It gives neither a computational procedure on arbitrary surcomplex inputs
nor a canonical parametrization of all class automorphisms. The much more
explicit decomposition in Section~\ref{saut:sec:decomposition} succeeds precisely
because it retains a valuation and a normal-form presentation.

\section{Involutions and genuinely different real forms}\label{saut:sec:realforms}

An \emph{involution} in this section means an automorphism of order exactly
two. A \emph{real form} of $\K$ means a real closed subfield $F\subset\K$
such that $\K=F(i)$. Its conjugation is the involution fixing $F$ and sending
$i$ to $-i$. The fields $F_\theta$ from \eqref{saut:eq:twistedconj} already give
many examples, all isomorphic to the distinguished $\No$.

We use the classical Artin--Schreier theorem: if an algebraically closed
field is a proper finite extension of a subfield, that subfield is real
closed, the characteristic is zero, and the extension has degree two.
A proof and the finite-automorphism application are given in \cite{Conrad}.
We apply these ordinary set-field theorems inside invariant set-sized
algebraically closed subfields, as follows.

\begin{theorem}[Finite subgroups and real forms]\label{saut:thm:torsion}
Every nontrivial finite subgroup of $\Aut(\K)$ has order two.
Every nonidentity finite-order automorphism is an involution $j$, its fixed
field $F_j$ is real closed, and $\K=F_j(i)$. In particular, $j(i)=-i$.
Consequently $\Aut(\K/\C)$ and the leading-term kernel $\U$ are torsion-free.
Within $G_{\No}$, the only nonidentity torsion element is $c$.
\end{theorem}
\begin{proof}
For a nontrivial finite indexed group $H$, choose for each nonidentity
element a moved point. Adjoin $i$ and close this finite set under $H$,
then take the relative algebraic closure in $\K$ of its generated field.
The result is a set-sized, $H$-stable algebraically closed field $E$,
and the restricted action is faithful. Artin's fixed-field theorem and
Artin--Schreier give $[E:E^H]=|H|=2$. Applying this to a finite cyclic
action shows that every nonidentity finite-order automorphism is an
involution.

For such an involution $j$, the same construction can include any prescribed
set of elements of $\Fix(j)$. Each $E^j$ is real closed, and these fields
form a directed family with union $F_j=\Fix(j)$; hence $F_j$ is real closed.
Because $i\in E$ and $E^j$ is real closed, $j(i)=-i$. For every $z\in\K$,
\[
 z=\frac{z+jz}{2}+i\frac{z-jz}{2i},
\]
and both displayed coordinates are fixed by $j$. Thus $\K=F_j(i)$.
An automorphism fixing $\C$ cannot be such an involution.
A member of $\U$ must fix $i$, since changing it to $-i$ changes its leading
coefficient, and is therefore also torsion-free.
Finally, an order-preserving bijection of a linear order has no nontrivial
finite order: a moved point has a strictly monotone finite orbit segment.
Use $G_{\No}\simeq\Aut(\No)\times C_2$.
\end{proof}

Two involutions are conjugate in $\Aut(\K)$ exactly when their fixed real
fields are isomorphic. One direction follows by restricting a conjugating
automorphism. Conversely an isomorphism $F_j\to F_k$ extends by
$x+iy\mapsto f(x)+if(y)$, and this extension conjugates $j$ to $k$.
Thus classifying involutions is a real-form classification problem, not
merely the observation that $i$ has two possible images.

\subsection{A real form not isomorphic to the surreal field}

The following construction shows that not all involutions of $\K$ are
conjugate to its distinguished conjugation.

\begin{theorem}\label{saut:thm:inequivalentrealform}
Assuming the class-theoretic framework above, $\K$ has a real form with a
countable cofinal subset. It is not isomorphic as a field to $\No$.
The associated involution is not conjugate to $c$.
\end{theorem}
\begin{proof}
Form the rational function field $\No(T)$ with $T$ transcendental, ordered
so that $T$ is larger than every surreal coefficient. The sign of a
nonzero rational function is determined by the ratio of its leading
coefficients. Let
\[
 H=\No(T)^{\rc}
\]
be an ordered real closure. To construct this in class foundations,
choose a set-like enumeration of $\No(T)$ and increasing set-sized subfields
$A_\alpha$ exhausting it. For each $\alpha$ choose a set-coded ordered real
closure $R_\alpha$ of $A_\alpha$, using the fixed global well-order.
For $\alpha\le\beta$, there is a unique ordered embedding
$R_\alpha\to R_\beta$ extending $A_\alpha\subseteq A_\beta$.
Uniqueness follows by preserving the order of the finitely many real roots
of each polynomial, and makes these embeddings compatible under composition.
Take the direct limit: represent an element by the least equivalent
stage--element pair in the set-like global well-order. Thus the elements
are sets, not proper-class equivalence classes. Every finite tuple occurs
in one stage, so the resulting class field is real closed; every element
is algebraic over the embedded $\No(T)$. Identifying that canonical image
with $\No(T)$ gives the required extension. This is a coherent set-stage
construction, not an application of a class Zorn lemma.

The powers $T,T^2,T^3,\ldots$ are cofinal in $\No(T)$. Indeed, each rational
function has finite degree and each nonzero coefficient and its reciprocal
is bounded in absolute value by $T$. Increasing an integer power of $T$
therefore dominates the function. The same powers remain cofinal in $H$:
if $h$ is algebraic, a monic polynomial for $h$ has finitely many
coefficients, and the usual algebraic root bound
\[
 |h|\le 1+\max_{k<n}|a_k|
 \quad\text{when }h^n+a_{n-1}h^{n-1}+\cdots+a_0=0
\]
is valid in any ordered field. Its right-hand side is bounded by a power
of $T$.

The field $H$ is proper-class-sized because it contains $\No$. Its
complexification $H(i)$ is algebraically closed of characteristic zero.
By \cref{saut:thm:acfhomogeneity}, it is isomorphic as a pure field to $\K$.
Transport $H$ and its conjugation across this isomorphism.
The transported real form has a countable cofinal subset. By contrast,
every set-sized subset of $\No$ is bounded above, so $\No$ has no countable
cofinal subset. A field isomorphism between real closed fields preserves
their unique orders; hence these real forms cannot be isomorphic.
The conjugacy criterion for involutions proves the final assertion.
\end{proof}

This separates two phenomena. The phase family gives distinct embedded
copies of the \emph{same} real form. Theorem~\ref{saut:thm:inequivalentrealform}
gives a different isomorphism type of real form. Neither result supplies a
classification of all real closed class fields whose complexification is
$\K$.

\section{Simplicity, the omega-map, and exponential structure}\label{saut:sec:exponential}

The large automorphism groups above forget much of what makes a surreal
number canonical. Adding back suitable structure can impose rigidity, but
different added structures impose different conditions.

\subsection{Simplicity restores the canonical surreal numbers}

Let $<_{s}$ denote proper-prefix simplicity on ordinal-length surreal sign
sequences. A bijection preserving both numerical order and simplicity is
identity. For completeness, fix the empty sign sequence. At a successor
stage, the two immediate simplicity extensions of a fixed sequence lie on
opposite numerical sides, so both are fixed. At a limit stage, the unique
sign sequence whose proper initial segments are the already fixed chain
is fixed. Transfinite induction proves the assertion.

Thus a member of $G_{\No}$ whose real restriction preserves simplicity is
either identity or conjugation; marking $i$ eliminates conjugation. This
is rigidity of an enriched structure, not rigidity of the real closed
field reduct. A transformation such as $S_2$ is allowed to change birthdays
and simplicity and therefore does not contradict it.

\subsection{The omega-map detects exact value preservation}

Suppose $\sigma\in\Aut(\No)$ commutes with $x\mapsto\omega^x$ and fixes
valuation values pointwise. Since $v(\omega^x)=-x$,
\[
 -\sigma(x)=v(\omega^{\sigma(x)})
 =v(\sigma(\omega^x))=v(\omega^x)=-x.
\]
Therefore $\sigma=\id$. The qualifier ``pointwise'' is indispensable:
all surreal field automorphisms preserve the natural valuation \emph{up to}
an ordered group automorphism, and that weaker fact does not give the last
equality. Without pointwise fixation, the displayed computation instead
says that the induced value-group action equals $\sigma$ itself on the
underlying surreal elements. That is a compatibility condition, not a
proof of triviality.

\subsection{Exponential faithfulness: the existing repository result}

The following finite argument is taken from the repository's
\emph{Growth-Scale Rigidity} manuscript \cite{RepoRigidity}. We give the
proof because it is short, powerful, and avoids exchanging an automorphism
with an infinite sum.

\begin{lemma}[Displacement cannot be globally bounded]\label{saut:lem:displacement}
Let $F$ be an ordered field and $\sigma:F\to F$ a nonidentity unital field
endomorphism. Then $\{\sigma(x)-x:x\in F\}$ is cofinal and coinitial in $F$.
\end{lemma}
\begin{proof}
Write $D(x)=\sigma(x)-x$. The correct multiplication identity is
\begin{equation}\label{saut:eq:displacement}
 D(ab)=\sigma(a)D(b)+bD(a).
\end{equation}
Choose $a$ with $d=D(a)\ne0$. Given $B>0$, put
\[
 b=\frac{2B(1+|\sigma(a)|)}{|d|}.
\]
If both $|D(b)|$ and $|D(ab)|$ were at most $B$, then
\[
 2B(1+|\sigma(a)|)=|bd|
 \le |D(ab)|+|\sigma(a)||D(b)|
 \le B(1+|\sigma(a)|),
\]
a contradiction. Thus absolute displacements have no upper bound.
Since $D(-x)=-D(x)$, both signs occur with arbitrarily large magnitude.
\end{proof}

\begin{theorem}[Faithful value-group action]\label{saut:thm:expfaithful}
Let $F$ be an ordered field with an increasing group isomorphism
$E:(F,+)\to(F_{>0},\cdot)$, and let $w:F^\times\to\Lambda$ be a nontrivial
convex valuation. If an automorphism $\sigma$ commutes with $E$ and fixes
all $w$-values, then $\sigma=\id$.
Hence the exponential automorphisms stabilizing the valuation ring act
faithfully on its value group.
\end{theorem}
\begin{proof}
The homomorphism
\[
 \nu:F\to\Lambda,\qquad \nu(x)=w(E(x))
\]
is additive, with image exactly $w(F^\times)$: for $a\ne0$ use
$x=E^{-1}(|a|)$ and $w(-1)=0$ to get $\nu(x)=w(a)$.
Thus it is onto the actual value group, even if $\Lambda$ was presented
as a larger ambient group. Nontriviality of $w$ makes
its kernel $H$ a proper convex additive
subgroup: $E$ is increasing and a convex valuation reverses the order on
positive elements. A proper convex additive subgroup is bounded above in
$F$; if $b>0$ lies outside it, then every $h\in H$ satisfies $|h|<b$.

Commutation and pointwise value fixation give, for every $x$,
\[
 \nu(\sigma x-x)
 =w(E(\sigma x))-w(E(x))
 =w(\sigma(E(x)))-w(E(x))=0.
\]
Thus all displacements lie in the bounded subgroup $H$.
Lemma~\ref{saut:lem:displacement} forces $\sigma=\id$.
Taking the kernel of the induced-action homomorphism proves faithfulness.
\end{proof}

For the usual surreal exponential, the kernel $H$ is exactly the finite
real-surreal elements. Indeed, bounded $x$ gives bounded $\exp x$ and
bounded reciprocal by comparison with $e^n$; conversely, if $\exp x$ and
its reciprocal are bounded, comparison with integers and application of
$\log$ bound $x$. Consequently,
\begin{equation}\label{saut:eq:expfaithful}
 \ker\bigl(\Gexp\longrightarrow\Gh\bigr)=\{\id\}.
\end{equation}
In particular, a nonidentity surreal $1$-automorphism cannot preserve the
total real exponential. This settles the precise mathematical alternative
posed in Question~5.4 of the inspected arXiv version of \cite{KKS}, by the
repository proof above; it is not a new result of this article.
Faithfulness does not say that the exponential automorphism group is
trivial. It says that no nonidentity member can hide all its motion below
the valuation scale.

\subsection{A simple obstruction to exponential lifting}

An ordered automorphism of the value group always has the Hahn lift
$M_{\id,\tau}$. It need not have \emph{any} lift preserving the total
surreal exponential. One elementary instance, also covered in
\cite{RepoRigidity}, is worth spelling out.

\begin{proposition}\label{saut:prop:rationalnonlift}
For $q\in\Q_{>0}$ with $q\ne1$, the value-group dilation $g\mapsto qg$
is not induced by any exponential automorphism of $\No$.
\end{proposition}
\begin{proof}
If $\sigma$ were such an automorphism, the homomorphism
$\nu(x)=v(\exp x)$ would satisfy
$\nu(\sigma x)=q\nu(x)=\nu(qx)$, using the $\Q$-linearity of the additive
map $\nu$. Thus
$\sigma(x)-qx$ would be finite for every $x$.
Choose a positive infinite $X$, and write
\[
 \sigma(X)=qX+e,\qquad \sigma(X^2)=qX^2+f
\]
with $e,f$ finite. Multiplicativity yields
\[
 q^2-q+\frac{2qe}{X}+\frac{e^2-f}{X^2}=0.
\]
The last two terms are infinitesimal. Taking standard part gives
$q^2-q=0$, contradicting $q>0$, $q\ne1$.
\end{proof}

Thus even the very simple field automorphism $S_2$ cannot be repaired by
composing with a value-fixing automorphism to obtain an exponential lift
of its value action.

\subsection{A global surcomplex operation that remembers real exponential}

There is no need to postulate a total complex logarithm. The operation
\begin{equation}\label{saut:eq:L}
 L:\K^\times\to\No\subset\K,
 \qquad L(z)=\log|z|
\end{equation}
is already globally defined, using only the usual real surreal logarithm.
Say that $\alpha$ preserves $L$ when
$\alpha(L(z))=L(\alpha(z))$ for every nonzero $z$.

\begin{theorem}[Logarithmic-modulus automorphisms]\label{saut:thm:L}
For $\K=\No(i)$,
\begin{equation}\label{saut:eq:LAUT}
 \AutL(\K)\simeq\Gexp\times C_2.
\end{equation}
Its action on the natural value group has kernel exactly $\{\id,c\}$.
Marking $i$ makes that action faithful.
\end{theorem}
\begin{proof}
The range of $L$ is all of $\No$, because $L(\exp x)=x$. A commuting
automorphism and its inverse therefore preserve $\No$. Write it as
$\widehat\sigma c^e$. For $r>0$ real-surreal,
\[
 \sigma(\log r)=\alpha(L(r))=L(\sigma r)=\log(\sigma r).
\]
Thus $\sigma$ commutes with logarithm and hence with exponential.
Conversely, if $\sigma$ is exponential, the norm equivariance
\eqref{saut:eq:equivariantnorm} proves that both lifts commute with $L$.
The kernel statement follows from \cref{saut:thm:expfaithful} and the fact that
conjugation fixes every valuation value.
\end{proof}

Theorem~\ref{saut:thm:L} is another result already present in
\cite{RepoRigidity}. Notice that $L$ is required to transform
\emph{equivariantly}: $L(\alpha z)=\alpha(Lz)$. Demanding instead
$L(\alpha z)=Lz$ pointwise would impose a stronger condition.

\subsection{Local exponentials and derivations are different structures}

For $\varepsilon\in\mm$, the Hahn series
\[
 \operatorname{Exp}(\varepsilon)=\sum_{n\ge0}\frac{\varepsilon^n}{n!},
 \qquad
 \operatorname{Log}(1+\varepsilon)
   =\sum_{n\ge1}\frac{(-1)^{n+1}\varepsilon^n}{n}
\]
are well-defined by positive-support summability. Any strongly
$\C$-linear element of $G_v$ commutes with these operations, simply by
applying strong summation and the field laws. This includes the phase
twists, the exponent dilations, and the constant-fixing flows above.
For finite $z=a+\varepsilon$, such an automorphism also respects the
finite exponential $e^a\operatorname{Exp}(\varepsilon)$.
None of this implies compatibility with the total real surreal exponential
on infinite arguments, as $S_2$ and \cref{saut:prop:rationalnonlift} demonstrate.

The Berarducci--Mantova construction equips $\No$ with a nonzero derivation
compatible with substantial exponential and transseries structure
\cite{BM}. Fix one such derivation $\partial$ and extend it to $\K$ by
\[
 \partial(x+iy)=\partial x+i\partial y.
\]
For real-axis-preserving automorphisms, a direct coordinate calculation gives
\[
 \{\alpha\in G_{\No}:\alpha\partial=\partial\alpha\}
 \simeq\Aut(\No,\partial)\times C_2.
\]
This is a reduction, not a classification of the differential group or of
all differential automorphisms without a marked real axis. In particular,
the scalar derivation $\partial$, formal differentiation of a power series,
and the fine derivative of a function are three distinct notions. The
zero-derivative automorphism in \cref{saut:thm:zeroderiv} concerns the third.

\section{Consequences for algebra, geometry, and symbolic computation}\label{saut:sec:applications}

Automorphisms transport identities, but a transported identity must include
all its coefficients and all its chosen structures. This elementary rule is
especially useful when applying surcomplex theory.

\subsection{Polynomial and finite-dimensional invariants}

If $p(X)=\sum_{k=0}^n a_kX^k$ and $\alpha$ is a field automorphism, put
$p^\alpha(X)=\sum_{k=0}^n\alpha(a_k)X^k$. Then
\[
 \alpha(p(z))=p^\alpha(\alpha z).
\]
Roots, root multiplicities, common-root relations, resultants,
discriminants, and finite matrix determinants therefore transform
algebraically. When all coefficients are fixed, the root multiset is
permuted. Without coefficient fixation, it is the roots of the
\emph{transformed} polynomial that are obtained.

For $\alpha\in G_v$, root separations satisfy
\[
 v(\alpha r_j-\alpha r_k)=\tau_\alpha(v(r_j-r_k)).
\]
Thus finite valuation-cluster trees retain their combinatorial incidence,
although the labels can be reindexed. For $\alpha\in\U$, every nonzero
separation even retains its leading term. These statements follow from
additivity and the valuation laws; they require no compactness or ordinary
analytic convergence.

For $\alpha\in G_{\No}$, Hermitian algebra is also transported, because
$\alpha$ commutes with conjugation. The identity
$A^*=\overline A^{\,T}$ transforms to $(\alpha A)^*=\alpha(A^*)$,
and positive real scalars remain positive. Under a phase twist not
preserving the axis, the same conclusion holds only after conjugation is
also transported to $j_\theta$. Keeping the original conjugation while
changing the scalars is not an invariant operation.

\subsection{Automorphism tests as checks on a proposed canonical definition}

Suppose a construction is claimed to depend only on the pure field
$(\K,+,\cdot)$ and on a specified set of scalar parameters. It must commute
with every automorphism fixing those parameters. Proposition~\ref{saut:prop:nondefinability}
then immediately rules out pure-field definitions of the distinguished
real axis and finite ring. A purported canonical modulus on the bare
field would similarly determine extra real-form information and cannot
be the chosen modulus without additional structure.

This does not make those constructions invalid. It identifies which data
must be declared: conjugation or a real form for Hermitian geometry,
a valuation for scales, monomials and a summation structure for normal-form
algorithms, and an exponential for logarithmic growth. An automorphism
counterexample is often the quickest way to detect an omitted hypothesis.

\subsection{What a symbolic implementation can realistically support}

A symbolic system can implement \eqref{saut:eq:monomial} exactly on a declared
Hahn workspace once it can evaluate the coefficient automorphism, exponent
map, and character, and certify that the exponent map is an ordered
additive bijection. For finite supports, the arithmetic checks reduce to
finite convolution. For symbolic infinite families, well-order and
finite-incidence certificates remain necessary.

The maps $S_a$, $D_\lambda$, and $P_\theta$ are particularly transparent on
normal-form expressions with accessible exponent coefficients. The
fixed-shift flow can be represented by its operator formula or by closed
expressions such as $t\mapsto t/(1-st)$. The coefficient-motion example
requires an abstract derivation on real constants and should not be
advertised as an algorithm on arbitrary numerical real inputs.
Likewise, the back-and-forth existence theorems are not executable
algorithms on unrestricted surcomplex names.

\section{Relation to the Surreal repository and a formalization plan}\label{saut:sec:repo}

This section is the source's comparison, as of its pin. It was checked for
this report and is accurate at the pin; Section~\ref{saut:sec:position}
adds the relations it did not draw, and
Appendix~\ref{saut:app:pinned} records what has changed since.

The inspection in this article uses the repository revision
\begin{center}
\repoSHA.
\end{center}
The root documentation, the report catalogue, the exponential-rigidity
manuscript and its audit, the normal-form bridge note, and the two specific
Lean modules below were read through the GitHub connector. This is a
selected-source inspection, not a claim to have searched every file or to
have rebuilt the project.

The inspected documentation distinguishes research arguments from checked
Lean coverage. In particular, the current bridge note reports an ordered
field equivalence for small normal forms, transfer of real closedness,
actual surcomplex algebraic closedness, and valuation/standard-part
preservation \cite{RepoBridge}. These are \emph{repository-reported}
implementation facts. This article supplies no independent build result
for them.

\subsection{Existing interfaces actually inspected}

\texttt{Surreal/Algebra/Complexify.lean} defines the quadratic algebra
$F[i]$, its imaginary unit, coordinate multiplication, conjugation,
\texttt{normSq}, and the identity \texttt{mul\_conj}. Its field construction
uses an ordered coefficient field and explicitly does not, by itself,
assert algebraic closedness \cite{RepoComplexify}.
These interfaces are sufficient foundations for a generic development of
\cref{saut:thm:axis,saut:thm:norm,saut:thm:circle}.

\texttt{Surreal/Algebra/FineDerivative.lean} defines
\texttt{FineHasDerivAt} by punctured-neighborhood difference quotients and
contains uniqueness, continuity, algebraic differentiation rules, and the
polynomial derivative theorem \cite{RepoFine}. It is an appropriate
interface for \cref{saut:thm:derivative}; the exponent-dilation example additionally
requires a proved normal-form automorphism and its exact valuation formula.

The documentation's statements about smallness must retain their universe
indices. A theorem that every lower-universe subset is closed and discrete
does not say that every subset available in a larger universe is discrete.
The class theory in this article and a universe-relative Lean realization
are compatible approaches, but they are not textually interchangeable.

\subsection{Suggested implementation order}

A useful order is to separate finite algebra from normal-form and global
existence arguments.

\begin{enumerate}
\item \textbf{Generic real-form algebra.} Construct the lift of a base-field
ring equivalence to $F[i]$. Prove uniqueness from its restriction and its
value on $i$, and identify the conjugation centralizer. Add exact norm
rigidity and the Cayley-circle criterion. These proofs need no Hahn theory.

\item \textbf{Explicit Hahn automorphisms.} Package coefficient maps,
ordered additive exponent equivalences, and characters in one construction.
Prove the composition law \eqref{saut:eq:monomiallaw} and valuation formula before
introducing the phase action. Keep the nested coefficient extraction
$\ell$ on surreal \emph{exponents} distinct from coefficient extraction
on a surcomplex input.

\item \textbf{Support-certified flows.} Reuse a summable-family interface
and prove the fixed-shift family is jointly summable. Establish inverse
formulas for \cref{saut:thm:coeffflow,saut:thm:shiftflow} through coefficientwise finite
identities. Merely entering a formal expression $\exp(D)$ is not a proof of
an automorphism.

\item \textbf{Topology and derivatives.} Prove the induced homeomorphism
from a valuation reindexing, the nonzero-derivative rigidity theorem for a
generic topological field, and finally $S_a'\equiv0$ for $a>1$ using
\eqref{saut:eq:dilationquot}. No mean-value theorem is needed or available for
that last conclusion.

\item \textbf{Global foundations separately.} Treat class back-and-forth,
real-form existence, and collection-level group factorizations in an
explicit foundational layer or in bounded universes with the appropriate
hypotheses. Do not silently replace a set-stage extension theorem by an
unrestricted type-level statement about every model.
\end{enumerate}

The names in this implementation order describe proposed work; they are
not assertions that corresponding new modules already exist. In particular,
no new Lean code or Lean-checking claim is included in the present package.

\section{What remains to be classified}\label{saut:sec:questions}

The central stabilizer question has a clean answer, and the valued subgroup
has an exact kernel--quotient decomposition. Neither result gives a fully
explicit description of the unrestricted automorphism collection.
The following are useful research targets, with their status kept distinct.

\paragraph{The image of exponential automorphisms on the value group.}
Faithfulness \eqref{saut:eq:expfaithful} determines the kernel, not the image.
An intrinsic necessary-and-sufficient criterion on an ordered additive
$\tau:\No\to\No$ for it to come from an exponential automorphism would
be much stronger. Positive rational dilations already show that arbitrary
Hahn lifting does not solve the problem. Any such criterion must use the
interaction between the value group and real exponential growth, not just
ordered-group structure. (This is Question~13.1, \lbl{q:image}, of the
collection's rigidity report, where it is also open.)

\paragraph{Omega-map compatibility.}
The nontrivial automorphisms preserving the omega-map, and extensions of
order automorphisms of its generalized epsilon fixed points, are posed in
Questions~5.6 and~5.7 of the inspected version of \cite{KKS}. They are not
resolved here. Exact value rigidity only describes one kernel and must not
be mistaken for a solution of those broader extension problems.

\paragraph{Effective descriptions inside the leading-term kernel.}
The abstract support-and-invertibility conditions for strong Hahn
substitutions are meaningful, but finite, compositional certificates for
larger symbolic families would be useful for computation and formalization.
The fixed-shift constructions are a tractable subclass, not an exhaustive
normal form for $\U$ or its strong part. This is a proposed computational
research direction, not a claim that the general Hahn literature lacks
abstract descriptions.

\paragraph{Real-form and continuity classifications.}
The phase family and the cofinality example show that real forms have both
embedding freedom and isomorphism-type freedom. A detailed classification
with additional valuation, summation, or simplicity conditions is a separate
problem. Similarly, this article proves that $G_v$ acts continuously but
does not identify every fine-continuous pure-field automorphism. These are
boundaries of the present account; they are not asserted to be newly
formulated or historically unresolved in every related setting.

\section{Source audit, proof status, and conclusions}\label{saut:sec:audit}

\subsection{Version-sensitive points}

The version-specific literature anchor is Kaplan--Krapp--Serra,
arXiv:2509.22374v3, submitted 23 April 2026, whose PDF is dated 27 April
2026. Its displayed formulas were checked against the PDF, not only the
HTML rendering. Two local cautions affect example construction. The
shift $g-1$ printed in Example~4.3 is not contracting for the convention
$v=\min\supp$; our construction uses a positive shift. The series
$\sum_n\omega^{-1/n}$ printed in Example~4.4 does not have decreasing
Conway support; \eqref{saut:eq:nonstronga} uses a valid support instead.
Our nonstrong example does not claim the stronger global leading-term
property asserted for that source's construction.

The semidirect actions in this article are specified by actual composition
formulas, especially \eqref{saut:eq:monomiallaw} and \eqref{saut:eq:skeletonlaw}.
A display listing factors should not be read as a componentwise direct
product when exponent reindexing acts on coefficient weights.

The repository's exponential-rigidity report is explicitly prior material
at the pinned revision. Its appended source audit records an earlier
catalogue/tree mismatch, illustrating why absence from a catalogue or
failure of code search is not evidence of absence from the repository.
The present article neither relies on that earlier absence claim nor
makes an exhaustive nonduplication claim.

\subsection{Evidence categories}

\begin{center}
\small
\begin{tabular}{@{}>{\RaggedRight\arraybackslash}p{0.26\textwidth}>{\RaggedRight\arraybackslash}p{0.65\textwidth}@{}}
\toprule
Category & Treatment in this article\\
\midrule
Classical prerequisites & Real closedness and normal forms of $\No$;
Newton--Puiseux; finite-group fixed-field theory and Artin--Schreier are
imported with references.\\[3pt]
Direct proofs & Real-axis, circle, and norm stabilizers; explicit Hahn maps
and flows; valued factorization; topology and derivative results;
class extension and real-form consequences.\\[3pt]
Repository prior work & Exponential displacement faithfulness, rational
non-lifting, and logarithmic-modulus reduction are attributed, with proofs
included for their role here.\\[3pt]
Implementation evidence & Selected source code and documentation were read;
no independent repository build and no new Lean certification.\\[3pt]
Historical novelty & Not certified. Derived statements and useful examples
are not automatically claimed to be first discoveries.\\
\bottomrule
\end{tabular}
\end{center}

\subsection{Non-claims}\label{saut:sec:nonclaims}

The source states its limitations where they arise, and each is kept in
place. The list collects them, with the place where each stands.
Items~1--28 are the source's; items~29--33 were added when the report
joined the collection.
\begin{enumerate}[label=\arabic*.,leftmargin=2.2em]
\item No complete classification of the automorphisms of $\K$, and in
particular none of the unrestricted ones, is claimed (title page, abstract,
Section~\ref{saut:sec:questions}; the delivered README).
\item Historical priority is not claimed or certified; derived statements
and useful examples are not automatically claimed to be first discoveries,
and no exhaustive nonduplication claim is made (title page,
Section~\ref{saut:sec:audit}; the delivered README).
\item No new Lean verification, no Lean code, and no independent build of
the repository are claimed; the implementation order describes proposed
work, not existing modules (title page, Section~\ref{saut:sec:repo},
Section~\ref{saut:sec:audit}).
\item The repository comparison is a selected-source inspection, not a
search of every file; the bridge-note facts are repository-reported, not
independently rebuilt (Section~\ref{saut:sec:repo}).
\item Exponential faithfulness, the rational non-lifting proposition and
the logarithmic-modulus reduction are prior work of the repository, not
results of this report; the settlement of Question~5.4 of \cite{KKS} is by
the repository's proof and ``is not a new result of this article''
(abstract, Sections~\ref{saut:sec:introduction}
and~\ref{saut:sec:exponential}, Section~\ref{saut:sec:audit}).
\item Faithfulness does not say that the exponential automorphism group is
trivial (Section~\ref{saut:sec:exponential}).
\item The article does not rely on the earlier absence claim in the
rigidity report's audit (Section~\ref{saut:sec:audit}).
\item $\Aut(\K)$ and the other groups are definable collections of class
maps, not class objects of NBG; a group isomorphism is a unique
factorization statement about individual automorphisms
(Section~\ref{saut:sec:introduction}).
\item Global homogeneity and the topology of the full field are not to be
transferred to an arbitrary fixed set-sized workspace
(Section~\ref{saut:sec:introduction}).
\item The four-layer decomposition \eqref{saut:eq:fourfactor} is not a
decomposition of all field automorphisms of $\K$, and its factor $\U$ is not
reduced to a finite list of coefficients (Section~\ref{saut:sec:decomposition}).
\item An arbitrary homomorphism $g\mapsto u_g$ into $1+\mm$ does not by itself
give a summable or surjective map; the fixed-shift flows illustrate the
automorphism--derivation connection without asserting that every formal
operator exponential is summable (Section~\ref{saut:sec:kernels}).
\item The Taylor motions $\Phi_d$ are generally noncomputable from
numerical names of real coefficients and are not to be advertised as
algorithms; the back-and-forth existence theorems are not executable
algorithms and give no canonical parametrization of the class
automorphisms (Sections~\ref{saut:sec:kernels}, \ref{saut:sec:bare}
and~\ref{saut:sec:applications}).
\item The maps of $G_{\No}$ are semilinearly compatible with the modulus;
calling them isometries would be misleading (Section~\ref{saut:sec:axis}).
\item Proposition~\ref{saut:prop:continuous} is an inclusion in the group of
fine-continuous automorphisms, not a classification of it, and no converse
is inferred from continuity; not every fine-continuous pure-field
automorphism is identified (Sections~\ref{saut:sec:topology}
and~\ref{saut:sec:questions}).
\item The non-strong witness of Section~\ref{saut:sec:topology} is not
claimed to fix every value or every leading term, nor to have the global
leading-term property asserted for the construction in \cite{KKS}; the
non-density conclusion concerns the pointwise-equality topology on
automorphisms only (Sections~\ref{saut:sec:topology}
and~\ref{saut:sec:audit}).
\item The choice of topology is part of each topological theorem; the
discreteness of set-sized subsets in the full fine topology says nothing
about the intrinsic valuation topology of a fixed Hahn workspace
(Section~\ref{saut:sec:topology}).
\item $S_2$ is not asserted to be representable by a locally convergent
power series; ``everywhere differentiable'' is not a synonym for analytic
rigidity, and a zero derivative does not imply local constancy
(Section~\ref{saut:sec:derivative}).
\item The scalar Berarducci--Mantova derivation, formal differentiation of a
power series and the fine derivative are three distinct notions; the
differential reduction is not a classification of the differential group
or of differential automorphisms without a marked real axis
(Section~\ref{saut:sec:exponential}).
\item The transcendence-basis description is not a classification, even for
a set-sized algebraically closed field (Section~\ref{saut:sec:bare}).
\item Neither the phase family nor
Theorem~\ref{saut:thm:inequivalentrealform} classifies the real closed
class fields whose complexification is $\K$; the conjugacy criterion of
Appendix~\ref{saut:app:conjugacy} is structural, not an effective decision
procedure (Section~\ref{saut:sec:realforms},
Appendix~\ref{saut:app:conjugacy}).
\item Simplicity rigidity is rigidity of an enriched structure, not of the
real closed field reduct (Section~\ref{saut:sec:exponential}).
\item Exact value rigidity for the omega-map needs pointwise value
fixation; it describes one kernel and is not a solution of the omega-map
extension problems, Questions~5.6 and~5.7 of \cite{KKS}, which are not
resolved (Sections~\ref{saut:sec:exponential}
and~\ref{saut:sec:questions}).
\item Compatibility with the local exponential and logarithm does not imply
compatibility with the total real exponential on infinite arguments
(Section~\ref{saut:sec:exponential}).
\item The image of exponential automorphisms in the value group is not
determined (Section~\ref{saut:sec:questions}).
\item Automorphism tests do not make the constructions they constrain
invalid; they identify data that must be declared
(Section~\ref{saut:sec:applications}).
\item Smallness statements keep their universe indices; the class theory
here and a universe-relative Lean realization are compatible but not
textually interchangeable (Section~\ref{saut:sec:repo}).
\item The fixed-shift constructions are a tractable subclass, not an
exhaustive normal form for $\U$; the research targets of
Section~\ref{saut:sec:questions} are not asserted to be newly formulated or
historically unresolved in every setting, nor is it claimed that the Hahn
literature lacks abstract descriptions (Section~\ref{saut:sec:questions}).
\item The delivered README reports symbolic finite checks of the norm
product and polarization identities, the Cayley inverse, and the
rational-flow composition and inverse identities, and says that such
checks are not proofs of infinite summability or of the full surreal
theory. No code or record of them was delivered, and nothing in this
report rests on them. The delivery redistributes no third-party papers or
font files.
\item \emph{(Added.)} The answer to Question~5.4 of \cite{KKS} rests on the
rigidity report's check of the pinned version of that paper. The source's
readings of its Examples~4.3 and~4.4 and of its Questions~5.6 and~5.7 are
the source's and were not rechecked for this report.
\item \emph{(Added.)} Whether $\Gexp$ is nontrivial, equivalently whether
$\{\id,c\}\subsetneq\AutL(\K)$, is not decided; nor is whether $S_a|_{\No}$
is exponential when $a$ is not rational (Section~\ref{saut:sec:conventions}).
\item \emph{(Added; status updated.)} No dedicated \texttt{saut:} mappings
are supplied. Existing generic Lean proofs cover ordered displacement,
exponential rigidity, logarithmic-modulus classification and the valuation
kernel, as recorded in
Section~\ref{saut:sec:position}; their actual surreal exponential
instantiation remains pending. The discreteness proof concerns
lower-universe-small subsets of the Lean carrier, not an unqualified
proper-class statement.
\item \emph{(Added.)} The relations of Section~\ref{saut:sec:position}
identify results by their content; no theorem is transferred between the
full fine topology and a workspace topology, or between the fine
derivative and another derivative notion.
\item \emph{(Added.)} Observation \eqref{saut:eq:Lchain} is an elementary
consequence of results printed here; it is not claimed as a result of the
source or of the rigidity report.
\end{enumerate}

\subsection{The main picture}

The correct starting point is the strict separation
\[
 \{\id,c\}=\Aut(\K/\No)
 \ \subsetneq\ G_{\No}
 \ \subsetneq\ G_v
 \ \subsetneq\ \Aut(\K).
\]
Preserving the real axis gives $\Aut(\No)\times C_2$; preserving the unit
circle gives the same group. Exact preservation of every modulus reduces
it to identity and conjugation. Pure-field automorphisms can move the
real axis, the finite ring, and even the isomorphism type of a chosen real
form under a change of conjugation.

The valuation organizes a large explicit subgroup but does not enforce
classical analytic behavior. Phase twists can move real monomials while
fixing constants and values; leading-term automorphisms can move ordinary
real constants by infinitesimals; fine-continuous maps need not preserve
Hahn sums; and exponent dilation can be a differentiable field
homeomorphism with derivative zero everywhere. Exponential structure
imposes a different and genuinely strong constraint: its value-group
action is faithful, even though the underlying valued-field kernel is
large.

An automorphism theorem for surcomplex numbers is therefore meaningful only
after its preserved structure, support convention, topology, and size
framework have been made explicit. With those choices in place, the theory
has both exact structural reductions and a rich supply of concrete examples.

\clearpage
\appendix
\section{A quick reference of explicit examples}\label{saut:sec:examplesummary}

In the table, ``values fixed'' means $v(\alpha z)=v(z)$ for every nonzero
$z$, not merely preservation up to a value-group automorphism.
The term ``constants'' means the embedded ordinary complex field.

\small
\begin{longtable}{@{}>{\RaggedRight\arraybackslash}p{0.20\textwidth}>{\RaggedRight\arraybackslash}p{0.35\textwidth}>{\RaggedRight\arraybackslash}p{0.35\textwidth}@{}}
\toprule
Map & Preserved data & Nontrivial feature\\
\midrule
\endfirsthead
\toprule
Map & Preserved data & Nontrivial feature\\
\midrule
\endhead
$c$ & Real axis pointwise; modulus; values; strong sums & $i\mapsto-i$;
not complex fine-differentiable\\[5pt]
$S_2$ & Constants; real axis; conjugation; strong sums & $v\mapsto2v$;
$\omega\mapsto\omega^2$; derivative zero\\[5pt]
$D_\lambda$, $\lambda\ne0$ & Constants; real axis; conjugation; values;
strong sums & Rescales leading coefficients by $e^{\lambda(g)}$\\[5pt]
$P_{\pi/2}$ & Constants; values; strong sums; fine topology & $t\mapsto it$;
moves the real axis and changes modulus\\[5pt]
$\Phi_d$, $d(b)=1$ & Monomials; leading terms; strong sums; real axis when
$d$ is real & $b\mapsto b+t^\delta$; ordinary constants need not be fixed\\[5pt]
$T_1$, $\lambda=\ell$, $\delta=1$ & Constants; real axis; conjugation;
leading terms; strong sums & $t\mapsto t/(1-t)$; not a total-real-exponential
automorphism\\[5pt]
Ordered extension of Section~\ref{saut:sec:topology} & Constants; real axis;
conjugation; fine topology & Fixes the summands of $a$ but sends
$a\mapsto a+t^\omega$\\[5pt]
Pure-field back-and-forth & Any prescribed set of parameters, with ordinary
$\C$ also fixed & Can send a positive infinite $x$ to $x^{-1}$ or $ix$;
no canonical normal-form action supplied\\
\bottomrule
\end{longtable}
\normalsize

\section{A conjugacy criterion for the standard real form}\label{saut:app:conjugacy}

The cofinality obstruction in \cref{saut:thm:inequivalentrealform} has a useful
positive counterpart. Say that an ordered class field has the
\emph{set-cut interpolation property} if for every pair of sets
$\mathcal L,\mathcal R$ with $\mathcal L<\mathcal R$, there exists $x$ with
$\mathcal L<x<\mathcal R$. Empty sides are permitted.

\begin{proposition}\label{saut:prop:conjugacycriterion}
An involution $j$ of $\K$ is conjugate to $c$ if and only if its fixed real
field has the set-cut interpolation property.
\end{proposition}
\begin{proof}
The property holds in $\No$ by Conway cuts and is preserved by ordered
isomorphisms. Conversely let $F=\Fix(j)$ have the property. It is a proper
class, since $\K=F(i)$, and it is real closed by
\cref{saut:thm:torsion}. Run the proof of \cref{saut:thm:orderedhomogeneity} between
$F$ and $\No$ rather than within $\No$. At every stage, the only special
existence condition used to extend the partial ordered-field isomorphism
is interpolation of a set-sized cut. It holds on both sides, so
back-and-forth gives $F\simeq\No$. The real-form conjugacy criterion in
Section~\ref{saut:sec:realforms} then gives $j\sim c$.
\end{proof}

This is a structural characterization, not an effective decision procedure
for a given class involution. It singles out precisely the conjugacy class
of the usual surreal real form among all real forms of the pure field.

\section{Provenance and the pinned comparison}\label{saut:app:provenance}
\RaggedRight

\subsection{Provenance of this report}

This report is built from \emph{one} manuscript: the archive
\texttt{surcomplex\_automorphisms}, manuscript~07 of batch~20 of the
collection, dated 22 September 2026 and pinned to repository commit
\repoSHA. It was placed in
commit \texttt{5fe7f8d} as a new report rather than merged into
\repo{surreal/exponential-automorphism-rigidity}: it continues that report's
section on $\No(i)$, but with plain field automorphisms instead of
exponential ones, and answers no question that report names. There was no
second source; nothing was merged and nothing was selected away. Every
result, proof, example, question and limitation of the manuscript is
printed here.

The delivery consisted of the article source, its PDF and a README; it
contained no code and no data. The source was renamed \texttt{article.tex}
and edited as listed below; the README was rewritten for the collection;
the delivered PDF is not shipped and is replaced by a build of this text.

The editorial changes are these.
\begin{enumerate}[label=(\arabic*)]
\item Every label carries the prefix \texttt{saut:}. The source's 79 labels
are kept, unchanged after the prefix, and seven were added
(\texttt{saut:sec:conventions}, \texttt{saut:sec:position},
\texttt{saut:eq:Lchain}, \texttt{saut:sec:nonclaims},
\texttt{saut:app:conjugacy}, \texttt{saut:app:provenance},
\texttt{saut:app:pinned}).
\item The title page gives the report's status and provenance, and the
abstract gains one sentence placing $\AutL(\K)$.
\item Sections~\ref{saut:sec:conventions} and~\ref{saut:sec:position} and
the consolidated non-claims of Section~\ref{saut:sec:nonclaims} are new.
Short pointers were added after Theorem~\ref{saut:thm:decomp}, at the start
of Section~\ref{saut:sec:repo}, and in the first question of
Section~\ref{saut:sec:questions}.
\item Six symbols were renamed so that each has one meaning
(Section~\ref{saut:sec:conventions}): $E_{\rho,\tau}\to M_{\rho,\tau}$,
$E_d\to\Phi_d$, the derivation $D\to\mathcal D$, the Puiseux field
$P\to\mathcal P$, the automorphism $\tau\to\psi$ of the non-density
argument, and cut sides $(L,R)\to(\mathcal L,\mathcal R)$; the source's
$\Aut(\K,L)$ is written $\AutL(\K)$.
\end{enumerate}
The placement edit preserved the source's mathematical statements.
It retained the results reproduced from the rigidity report (its
Lemma~3.1, the faithfulness of Corollary~4.2, the content of
Theorem~5.1 and Corollary~5.2, Theorem~8.4, and Theorem~9.2 with its
kernel clause). They are printed here once, with the
source's proofs and attribution, rather than replaced by pointers, because
they are part of one manuscript's argument; Section~\ref{saut:sec:position}
names each counterpart. The letter $E$ was freed for the exponential,
following the rigidity report, rather than renaming that report's symbol.
The title is the source's; the directory is named
\texttt{surcomplex-field-automorphisms} to say that the automorphisms are
field automorphisms.

The later proof review clarifies rational cuts and strong coefficient action.
It expands the class and support arguments and excludes every derivative
value in $\K$, including infinite values. The finite-observation
non-density proof is extended to all surcomplex field automorphisms.
Current generic Lean coverage is separated from the pinned inspection.
See \repo{REVIEW.md} for the changes and validation.
The added arguments are not attributed to the original delivery.

\subsection{Pinned repository scope}\label{saut:app:pinned}

Section~\ref{saut:sec:repo} is the source's comparison at its pin
\texttt{dcf8666}, an ancestor of the placement commit \texttt{5fe7f8d}
(sixteen commits earlier). Checked against the repository for this report:
\begin{itemize}
\item The descriptions of \nolinkurl{Surreal/Algebra/Complexify.lean},
\nolinkurl{Surreal/Algebra/FineDerivative.lean}, the normal-form bridge note,
and the rigidity report with its appended audit correction are accurate at
the pin. None of these files, nor the root \texttt{README.md}, changed
between the pin and the placement; the report catalogue
\nolinkurl{docs/README.md} changed only by listing reports of another batch
and extending other entries.
\item The bibliography cites the rigidity report by its article title,
\emph{Growth-Scale Rigidity for Surreal and Surcomplex Automorphisms}; the
collection's README for the same directory is headed ``Cofinal displacement
and the rigidity of exponential automorphisms''. Both name
\repo{surreal/exponential-automorphism-rigidity}.
\item The source did not relate its results to
\lbl{found:thm:discrete} and its Lean proof, to the collection's
zero-derivative witnesses, to the character twists of
\lbl{spec:lem:finiteindex}, to the Cayley parametrization of the
trigonometry report, or to the Lean modules
\nolinkurl{Surreal/Surcomplex/Basic.lean} and
\nolinkurl{Surreal/HahnSeries/ComplexNumbers.lean}, all present at the pin.
The hidden-negative-directions and bounded-support reports and
\nolinkurl{Surreal/Algebra/SigmaDerivation.lean} were added after the pin
and could not have been seen. Section~\ref{saut:sec:position} supplies
these relations.
\item The README of the rigidity report states that nothing there
proves a theorem about the automorphism group of $\No$ that any other
report also proves. That was accurate before this report was placed; this report
reprints several of its results (Section~\ref{saut:sec:position}).
\end{itemize}

\phantomsection
\addcontentsline{toc}{section}{References}
\begin{thebibliography}{99}
\small\RaggedRight
\bibitem{Conway}
J.~H.~Conway,
\emph{On Numbers and Games}, Academic Press, 1976.

\bibitem{Gonshor}
H.~Gonshor,
\emph{An Introduction to the Theory of Surreal Numbers},
London Mathematical Society Lecture Note Series 110,
Cambridge University Press, 1986.

\bibitem{KS}
S.~Kuhlmann and M.~Serra,
\emph{The automorphism group of a valued field of generalised formal power series},
Journal of Algebra \textbf{605} (2022), 339--376.
\href{https://doi.org/10.1016/j.jalgebra.2022.04.023}{doi:10.1016/j.jalgebra.2022.04.023}.
Preprint version inspected: \href{https://arxiv.org/abs/2107.03362v3}{arXiv:2107.03362v3},
11 April 2022.

\bibitem{KSG}
S.~Kuhlmann and M.~Serra,
\emph{Automorphisms of valued Hahn groups},
\href{https://arxiv.org/abs/2302.06290}{arXiv:2302.06290}, 2023.

\bibitem{KKS}
E.~Kaplan, L.~S.~Krapp, and M.~Serra,
\emph{Decomposing the automorphism group of the surreal numbers},
\href{https://arxiv.org/abs/2509.22374v3}{arXiv:2509.22374v3},
23 April 2026; PDF dated 27 April 2026.
The version-specific PDF, including pp.~8--10, was inspected.

\bibitem{BKKS}
V.~Bagayoko, L.~S.~Krapp, S.~Kuhlmann, D.~Panazzolo, and M.~Serra,
\emph{Automorphisms and derivations on algebras endowed with formal infinite sums},
\href{https://arxiv.org/abs/2403.05827v2}{arXiv:2403.05827v2},
22 September 2025.

\bibitem{BM}
A.~Berarducci and V.~Mantova,
\emph{Surreal numbers, derivations and transseries},
Journal of the European Mathematical Society \textbf{20} (2018), 339--390.
\href{https://doi.org/10.4171/JEMS/769}{doi:10.4171/JEMS/769}.
\href{https://arxiv.org/abs/1503.00315v3}{arXiv:1503.00315v3}.

\bibitem{Conrad}
K.~Conrad,
\emph{The Artin--Schreier Theorem}, author expository notes,
Sections 3--4, especially Theorem 3.1 and the finite-group application.
\url{https://kconrad.math.uconn.edu/blurbs/galoistheory/artinschreier.pdf}.
Accessed 22 September 2026.

\bibitem{Repo}
V.~Reshetnikov and contributors,
\emph{Surreal}, GitHub repository, snapshot
\repoSHA, inspected 22 September 2026.
\url{https://github.com/VladimirReshetnikov/Surreal}.
Root \texttt{README.md} and \texttt{docs/README.md} inspected.

\bibitem{RepoRigidity}
\emph{Growth-Scale Rigidity for Surreal and Surcomplex Automorphisms},
AI-assisted repository research manuscript, dated 21 September 2026;
\path{docs/surreal/exponential-automorphism-rigidity/article.tex},
at the snapshot in \cite{Repo}.
The accompanying \path{05-exponential-valuation-SOURCE_AUDIT.md}, including
its appended correction, was also inspected.
\href{https://github.com/VladimirReshetnikov/Surreal/blob/dcf86662b574988e75d08d515d3a2a5ba7bbc0d7/docs/surreal/exponential-automorphism-rigidity/article.tex}{Pinned manuscript source}.

\bibitem{RepoBridge}
\emph{Next steps for the infinite normal-form bridge},
\texttt{docs/NORMAL\_FORM\_BRIDGE.md}, snapshot in \cite{Repo},
opening assessment and implementation notes inspected.
\href{https://github.com/VladimirReshetnikov/Surreal/blob/dcf86662b574988e75d08d515d3a2a5ba7bbc0d7/docs/NORMAL_FORM_BRIDGE.md}{Pinned bridge note}.

\bibitem{RepoComplexify}
\nolinkurl{Surreal/Algebra/Complexify.lean}, snapshot in \cite{Repo}.
\href{https://github.com/VladimirReshetnikov/Surreal/blob/dcf86662b574988e75d08d515d3a2a5ba7bbc0d7/Surreal/Algebra/Complexify.lean}{Pinned source file}.

\bibitem{RepoFine}
\nolinkurl{Surreal/Algebra/FineDerivative.lean}, snapshot in \cite{Repo}.
\href{https://github.com/VladimirReshetnikov/Surreal/blob/dcf86662b574988e75d08d515d3a2a5ba7bbc0d7/Surreal/Algebra/FineDerivative.lean}{Pinned source file}.
\bibitem{ParanVo}
E.~Paran and T.~N.~Vo,
\emph{A skew Newton--Puiseux Theorem}, introduction (classical
Newton--Puiseux theorem), 2023.
\url{https://arxiv.org/abs/2311.17544}.
\end{thebibliography}
\end{document}
